% =============================================================================
% Neurogaze - AI Product Paper (versi final semifinal)
% Datathon RISTEK Fasilkom UI 2026, University Track
%
% Kompilasi:  pdflatex paper_final && pdflatex paper_final && pdflatex paper_final
%
% Gaya NeurIPS 2015. Berkas nips15submit_e.sty pada folder ini adalah
% reimplementasi spesifikasi resmi agar dokumen dapat dikompilasi di lingkungan
% tanpa berkas gaya asli. Bila berkas resmi tersedia, letakkan di folder ini dan
% berkas ini akan memakainya lebih dulu.
%
% Aturan redaksional yang mengikat seluruh naskah:
%   1. Setiap angka wajib memiliki artefak sumber (notebook, JSON, atau log).
%   2. Tidak ada klaim diagnosis, tidak ada klaim performa klinis pada balita.
%   3. Target pengembangan ditulis sebagai target, bukan sebagai hasil.
% =============================================================================

\documentclass{article}
\usepackage{nips15submit_e,times}
\usepackage[T1]{fontenc}
\usepackage[utf8]{inputenc}
\IfFileExists{bahasa.ldf}{\usepackage[bahasa]{babel}}{}
\usepackage{booktabs}
\usepackage{amssymb}
\usepackage{graphicx}
\usepackage{enumitem}
\usepackage{microtype}
\usepackage{hyperref}
\usepackage{url}

\hypersetup{colorlinks=true, linkcolor=black, citecolor=black, urlcolor=black}
\graphicspath{{../../figs/}}
\setlist{nosep, leftmargin=1.2em}

% Nama gambar, tabel, abstrak, dan daftar pustaka dalam bahasa Indonesia,
% tidak bergantung pada ketersediaan babel.
\renewcommand{\figurename}{Gambar}
\renewcommand{\tablename}{Tabel}
\renewcommand{\refname}{Referensi}
\renewcommand{\abstractname}{Abstrak}

% URL panjang boleh dipatahkan agar tidak melampaui blok teks.
\def\UrlBreaks{\do\/\do-\do.\do:\do_}
\Urlmuskip=0mu plus 3mu

\title{Neurogaze: Triase Skrining Dini Autisme Berbasis\\Pola Tatapan\\untuk Layanan Posyandu}

% Blok penulis disusun sebagai satu tabular tiga kolom agar ketiga nama tetap
% berada pada satu baris, sama seperti concept paper babak penyisihan.
\author{
\setlength{\tabcolsep}{3pt}
\begin{tabular}{ccc}
Timotius C. P. Sihombing & Raphael Angelo A. Purnama & Aqila Malfa Zahira \\[2pt]
{\normalfont\fontsize{8.5}{10}\selectfont\texttt{timotiuspatrick@gmail.com}} &
{\normalfont\fontsize{8.5}{10}\selectfont\texttt{raphael.adikara@gmail.com}} &
{\normalfont\fontsize{8.5}{10}\selectfont\texttt{aqilamalfa@gmail.com}}
\end{tabular}
}

\nipsfinalcopy

\setlength{\textfloatsep}{9pt plus 2pt minus 2pt}
\setlength{\floatsep}{8pt plus 2pt minus 2pt}
\setlength{\intextsep}{8pt plus 2pt minus 2pt}

% Dua diagram alur berukuran besar boleh berbagi satu halaman dengan teks;
% tanpa pelonggaran ini LaTeX melempar keduanya ke halaman terpisah dan
% menyisakan ruang kosong yang lebar.
\renewcommand{\topfraction}{0.92}
\renewcommand{\bottomfraction}{0.7}
\renewcommand{\textfraction}{0.06}
\renewcommand{\floatpagefraction}{0.85}
\setcounter{topnumber}{3}
\setcounter{totalnumber}{4}
\begin{document}

% ------------------------------------------------------------- halaman judul
\thispagestyle{empty}
\vspace*{1.1in}
\maketitle

\begin{center}
Fakultas Teknologi Maju dan Multidisiplin\\
Universitas Airlangga
\end{center}

\vfill

\begin{center}
\small AI Product Paper $\bullet$ Datathon RISTEK Fasilkom UI 2026 $\bullet$ Semifinal Round, University Track
\end{center}

\newpage
\setcounter{page}{1}

\begin{abstract}
Autisme paling responsif terhadap intervensi pada usia di bawah tiga tahun, tetapi tinjauan
lintas negara melaporkan diagnosis rata-rata pada usia 56 bulan, sekitar 32 bulan setelah
keresahan pertama orang tua; angka ini bukan estimasi Indonesia. Instrumen lini pertama
punya sensitivitas terbatas, sementara asesmen objektif terpusat di fasilitas mahal. Makalah
ini memperkenalkan Neurogaze, prototipe PWA luring untuk riset pengukuran atensi pada tablet
Android biasa. Kamera
mengestimasi pola pandangan selama stimulus sosial dan menahan sesi yang gagal memenuhi
syarat kualitas. Keputusannya berlapis tiga. GeoPref 69 persen adalah satu-satunya kandidat
pemicu rujukan dengan titik operasi terbit (1.863 anak usia 12--48 bulan, spesifisitas 98
persen, sensitivitas 17 persen), tetapi rujukan otomatis ditahan karena klip 16,75 detik
yang dikirim tidak mereplikasi protokol penuh tempat titik operasi itu diturunkan; enam indeks perilaku
dilaporkan deskriptif tanpa skor gabungan; dan lapis gabungan belum dibangun karena bobotnya
menuntut balita berlabel. Regresi logistik terkalibrasi yang dievaluasi di makalah ini tidak
memutuskan apa pun di produk dan dikurung pada panel riset berpenapis out-of-distribution.
Kami mengevaluasi bukti konsep domain sumber pada 547 scanpath dari 54 anak usia sekolah di
Prancis, direkam dengan eye-tracker 250~Hz, dengan pemisahan yang dikelompokkan per
partisipan tetapi tanpa uji eksternal. Regresi logistik 13 fitur geometri mencapai ROC-AUC
tingkat anak 0,823 (95\% CI 0,697-0,929) pada data tertahan, sedangkan CNN EfficientNetB0
pada citra yang sama mencapai 0,882 (0,774-0,968) tetapi tidak dapat dijalankan pada
perangkat sasaran. Pada proxy degradasi gabungan 30~Hz, derau 1,5$^\circ$, dropout 20\%, dan
jitter pose 2$^\circ$, fitur geometri bertahan pada AUC 0,683 dibandingkan 0,605 untuk 19
fitur yang bergantung pada encoding kinematik eye-tracker. Kontrol negatif memperlihatkan
taruhannya: mengganti unit pengelompokan split membuat 41 dari 54 anak muncul di kedua sisi.
Audit dataset wajah 2.940 citra menemukan enam dari enam metadata tata kelola tidak tersedia;
CNN di atasnya mencapai AUC 0,932, dan justru karena itu dataset dikarantina. Di luar
validasi algoritmik, dua gerbang instrumen sudah dituntaskan: Gate~A lulus pada 100 sesi dari
25 peserta dewasa lintas tiga perangkat Android dengan galat kalibrasi median 2,2$^\circ$, completion rate
94 persen, dan frame wajah-iris valid 96,4 persen; Gate~B lulus pada 30 perbandingan browser
simultan terhadap WebGazer.js 3.5.3, dengan 27 pasangan siap, galat median ternormalisasi
0,040997 dan agreement AOI rata-rata 99,7118 persen hasil rekomputasi dari koordinat mentah.
Blok target diketahui Gate~A menghasilkan konversi sudut lama 2,36$^\circ$ median dan
3,58$^\circ$ persentil 90 pada 94 sesi lulus, tetapi jarak pandang tidak direkam sehingga
angka itu bukan akurasi absolut eksak atau uji tanding terhadap WebGazer. Kontrol positif
dewasa menjadi bukti end-to-end utama: 12 peserta, 23 sesi direkam, 15 lulus mutu; 9/11
sesi menonton biasa dan 6/12 sesi pola diproduksi dapat dipakai, dengan aturan demonstrasi
menyala pada 0/9 dan 4/6 sesi. Hasil ini adalah manipulation check, bukan sensitivitas,
spesifisitas, ASD, atau status klinis peserta; mode demonstrasi tidak mengeluarkan rujukan.
Gate~C dan Gate~D tetap terbuka: belum ada balita, uji kader, validasi Indonesia,
wawancara praktisi, atau linkage klinis.
\end{abstract}

% =============================================================================
\section{Pendahuluan}

\subsection{Latar belakang}

Gangguan spektrum autisme menetap seumur hidup, tetapi lintasan perkembangannya dapat diubah
bila intervensi dimulai cukup awal. Uji acak terkendali Early Start Denver Model pada anak 18
sampai 30 bulan mencatat kenaikan skor kognitif 17,6 poin setelah dua tahun, dibandingkan 7,0
poin pada kelompok pembanding, dengan perubahan terbesar pada anak termuda~\cite{dawson2010}.
Persoalannya, sistem yang ada tidak menemukan anak pada rentang usia itu. Kajian lintas
negara melaporkan usia rata-rata diagnosis 55,97 bulan, padahal keresahan pertama orang tua
muncul pada 23,64 bulan~\cite{vanthof2021}: jeda 32,33 bulan yang jatuh persis pada periode
ketika intervensi paling berdampak. Ketimpangannya lebih tajam di negara berpenghasilan
rendah dan menengah, tempat 40 persen individu autistik baru terdiagnosis pada usia sepuluh
tahun atau lebih, dibandingkan 15 persen di negara berpenghasilan tinggi. Perbedaan itu
muncul di ekor distribusi, dan anak-anak di ekor itulah sasaran produk kami.

Instrumen skrining autisme yang paling luas dipakai adalah Modified Checklist for Autism in
Toddlers (M-CHAT), kuesioner yang diisi orang tua. Validasi pada layanan primer melaporkan
sensitivitas 49,2 persen dengan spesifisitas 94,2 persen dan nilai prediksi positif 28,9
persen~\cite{harper2026}; studi lain pada populasi risiko rendah mencatat nilai prediksi
positif 36 persen~\cite{sturner2016}. Angkanya berbeda antar validasi, tetapi arahnya
konsisten: sebagian anak autistik lolos dari penyaringan. Kinerjanya juga bervariasi menurut
jenis kelamin, usia gestasi, dan bahasa anak, sehingga perlu dilengkapi sinyal objektif yang
tidak bergantung pada ingatan orang tua.

Indonesia belum punya data prevalensi autisme nasional. Angka yang beredar berasal dari
ekstrapolasi negara lain atau studi regional yang metodenya tidak
sebanding~\cite{fathmawati2023}, sementara skrining M-CHAT-R di klinik Jakarta menemukan 3,52
persen anak 16 sampai 30 bulan berisiko tinggi~\cite{aditya2021}. Kapasitas asesmen spesialis
juga terbatas, sehingga menaikkan sensitivitas skrining tanpa memperbaiki kapasitas triase
hanya memindahkan masalah menjadi antrean rujukan.

Ada satu aset yang sering diabaikan. Indonesia memiliki 338.881 Posyandu dengan 1.059.466
kader aktif, dan program Stimulasi, Deteksi, dan Intervensi Dini Tumbuh Kembang (SDIDTK)
sudah menjadi program pokok Puskesmas yang menjangkau hampir setiap desa setiap
bulan~\cite{kemenkopmk2023}. Yang terbatas kapasitas deteksinya, bukan jangkauannya. Studi
pemberdayaan kader menemukan pelaksanaan SDIDTK cenderung berhenti pada penimbangan dan
pengukuran tinggi badan karena kader belum terlatih mendeteksi penyimpangan tumbuh
kembang~\cite{mardhika2022}. Jaringan penyalurnya sudah ada; kemampuan mengukurnya yang
belum.

\subsection{Solusi yang diusulkan dan kontribusi}

Kami mengusulkan Neurogaze (\emph{Neuro}divergent + \emph{gaze}), alat triase skrining
berbasis pola tatapan yang dijalankan kader pada tablet Android biasa. Produk ini tidak
menggantikan asesmen klinis dan tidak menyaingi ketepatan alat laboratorium. Perannya
melengkapi ceklis dengan pengukuran objektif untuk membantu keputusan rujukan pada titik
layanan yang sudah menjangkau seluruh desa. Dibandingkan concept paper babak penyisihan,
masalah dan produknya sama; yang bertambah adalah implementasi yang berjalan, validasi
algoritmik yang tuntas, dua gerbang instrumen yang sudah lulus pengujian, dan lingkup bukti
yang dinyatakan tegas.

Kontribusi makalah ini ada delapan:

\begin{enumerate}[label=(\roman*)]
\item rancangan triase untuk keterbatasan nyata layanan primer Indonesia (perangkat murah,
      operator bukan klinisi, konektivitas tidak menentu), diwujudkan sebagai PWA luring yang
      dideploy publik;
\item arsitektur keputusan berlapis yang memisahkan kewenangan dari perhitungan: satu pemicu
      rujukan yang ambangnya diambil apa adanya dari literatur, satu lapis indeks yang
      deskriptif dan dilarang digabung oleh tipe datanya sendiri, dan satu lapis gabungan
      yang sengaja belum dibangun (Bagian~\ref{sec:lapisan});
\item validasi algoritmik yang tuntas pada 54 anak dengan pemisahan berkelompok per
      partisipan, disertai kontrol negatif yang memverifikasi splitnya tahan bocor;
\item validasi instrumen yang tuntas pada dua gerbang: Gate~A pada 100 sesi dewasa lintas
      perangkat, pencahayaan, dan penggunaan kacamata, serta Gate~B pada 30 perbandingan
      browser simultan terhadap WebGazer.js;
\item perbandingan terkendali antara representasi yang akurat dan yang dapat dijalankan,
      yaitu CNN scanpath melawan 13 fitur geometri, beserta alasan model ber-AUC lebih tinggi
      tidak kami pakai;
\item studi degradasi yang memetakan penurunan performa ketika sinyal eye-tracker 250~Hz
      didekati kamera konsumen, dan yang mendasari penguncian set fitur produksi;
\item audit tata kelola yang menolak sebuah dataset wajah publik meski CNN di atasnya
      mencapai AUC 0,9324 (Lampiran~D);
\item pelaporan bukti berjenjang (Gate~A-D) yang \emph{fail-closed}, sehingga produk secara
      struktural tidak punya jalur untuk menampilkan angka validasi yang datanya belum
      terkumpul.
\end{enumerate}

\subsection{Batasan penelitian}
\label{sec:batasan}

Lima batasan berlaku pada seluruh hasil di makalah ini. Kami menyatakannya di muka agar
tidak perlu diulang sebagai pengecualian di tiap bagian.

\emph{Pertama, lingkup validasinya algoritmik dan instrumental.} Kami menguji tiga hal.
Apakah representasi dan model memuat sinyal yang dapat digeneralisasi ke subjek yang belum
pernah dilihat. Apakah kamera tablet cukup presisi untuk mengukur arah pandangan (Gate~A).
Dan apakah aliran pengukurannya sejalan dengan WebGazer.js sebagai referensi proyek (Gate~B). Ketiganya sudah
selesai. Validasi klinis pada kohort sasaran belum, sebab menuntut pengumpulan data primer
pada balita bersama mitra berizin etik. Gate~C dan Gate~D karena itu masih terbuka,
dan tidak ada klaim akurasi diagnosis di makalah ini.

\emph{Kedua, data anchor-nya berpindah domain dari sasaran produk.} Dataset Carette berisi
anak berusia rata-rata 7,88 tahun dari populasi Prancis yang direkam eye-tracker laboratorium
250~Hz, sedangkan sasaran produk anak 16-30 bulan di Indonesia pada kamera tablet 30~fps.
Studi degradasi menjawab sebagian dari perpindahan instrumennya; perpindahan usia dan
populasi hanya dapat dijawab kohort prospektif.

\emph{Ketiga, ukuran sampelnya kecil untuk standar epidemiologi.} Lima puluh empat anak
menghasilkan interval kepercayaan lebar dan titik kerja yang tidak stabil, sehingga kami
memperlakukan AUC sebagai hasil dan ambang operasional sebagai parameter per lokasi. Performa
per subkelompok tidak dapat dilaporkan karena dataset anchor tidak menyediakan demografi.

\emph{Keempat, keluarannya rekomendasi rujukan, bukan diagnosis.} Sistem tidak menyatakan
seorang anak autistik, melainkan bahwa pola atensinya di luar rentang tipikal sehingga layak
dinilai tenaga yang berwenang.

\emph{Kelima, belum ada satu pun balita dalam bukti apa pun di makalah ini, dan tidak satu
pun instrumen yang dipakai divalidasi di Indonesia.} Gate~A dan Gate~B seluruhnya pada
peserta dewasa, ambang Lapis~A berasal dari kohort Amerika Serikat, dan referensi Gate~B
divalidasi pada laboratorium di luar Indonesia. Batasan ini tidak dapat ditutup oleh
pekerjaan teknis; ia menuntut kohort prospektif lokal, dan konsekuensinya dirinci pada
Bagian~\ref{sec:tatakelola}.

% =============================================================================
\section{Tinjauan Pustaka}

\subsection{Atensi visual sebagai penanda perilaku}

Kelainan pola pandang adalah penanda perilaku autisme yang paling konsisten. Studi
longitudinal pada bayi usia 2 sampai 6 bulan menemukan penurunan fiksasi pada mata yang
tidak terjadi pada bayi yang kemudian berkembang tipikal~\cite{jones2013}, dan anak dengan
autisme memandang pola geometris jauh lebih lama dibandingkan adegan
sosial~\cite{pierce2011}. Sinyalnya muncul sebelum kemampuan berbahasa berkembang, sehingga
dapat diukur pada usia ketika instrumen berbasis wawancara belum dapat diterapkan.

Preferensi geometrik itulah yang kemudian divalidasi pada skala yang cukup besar untuk
melahirkan ambang. Wen dkk.~\cite{wen2022} menguji ambang 69 persen waktu fiksasi pada pola
geometrik terhadap 1.863 anak usia 12-48 bulan dan melaporkan sensitivitas 17 persen,
spesifisitas 98 persen, PPV 81 persen, dan NPV 65 persen. Bentuk angka itu penting dan
menentukan seluruh rancangan produk kami: ia menandai subtipe, bukan populasi ASD secara
keseluruhan, sehingga alat yang dibangun di atasnya hanya sah dipakai untuk menegakkan
kecurigaan dan tidak pernah untuk menyingkirkannya.

\subsection{Representasi scanpath sebagai citra}

Carette dkk.~\cite{carette2019} mengusulkan dasar teknis pekerjaan kami: mengubah rekaman
eye-tracking menjadi citra, dengan lintasan pandangan sebagai garis dan dinamikanya
dikodekan pada kanal warna, yaitu kecepatan, akselerasi, dan jerk, sehingga prediksi menjadi
klasifikasi citra. Mereka melaporkan AUC di atas 0,9 dengan jaringan saraf, sementara model
non-neural pada makalah yang sama hanya sekitar 0,7. Konsekuensinya jarang disebut: kanal
warna itu hanya bermakna pada laju sampling asalnya, sehingga fitur apa pun yang membacanya
ikut terikat pada eye-tracker 250~Hz.

\subsection{Praproses, fitur, model, dan protokol evaluasi}
\label{sec:tp-ringkas}

Empat keputusan metodologis di bawah menentukan angka mana yang berarti, dan rinciannya kami
pindahkan ke Lampiran~G agar dapat direplikasi tanpa memperpanjang badan makalah.

\emph{Praproses.} Tujuh tahap berjalan seluruhnya di perangkat: deteksi wajah dan iris lewat
MediaPipe Face Landmarker, kalibrasi affine ke koordinat layar, penapisan sampel di luar
margin, segmentasi deret pada jeda di atas 180~ms, penghalusan eksponensial dengan resampling
50~ms, pemetaan area minat, dan rasterisasi $640 \times 480$ mengikuti konvensi Carette.

\emph{Dua kelompok fitur.} Sembilan belas fitur terbelah menurut satu hal, yaitu apakah fitur
itu membaca kanal warna. Tiga belas fitur geometri hanya membaca posisi piksel bertinta
(cakupan, dispersi, pusat perhatian, jarak radial, entropi grid $8 \times 8$), sedangkan enam
fitur kinematik membaca kecepatan, akselerasi, dan jerk yang terkode pada kanal RGB. Turunan
tinggi menguatkan derau pada sampling rendah, jadi kelayakan kelompok kedua di kamera 30~fps
harus diputuskan lewat pengukuran, bukan asumsi.

\emph{Keluarga model.} Lima keluarga model tabular diuji pada kedua set fitur, yaitu regresi
logistik yang kami tetapkan sebagai model utama sebelum melihat hasil, SVM kernel RBF,
Gaussian Naive Bayes, Random Forest, dan Gradient Boosting, ditambah CNN
EfficientNetB0~\cite{tan2019} pada citra scanpath sebagai pembanding representasi.

\emph{Kalibrasi dan evaluasi.} Skor mentah dikalibrasi Platt pada lipatan validasi saja, dan
dua titik kerja dilaporkan berdampingan, yaitu ambang target sensitivitas dan indeks Youden.
Evaluasinya memakai GroupKFold bersarang dengan \texttt{participant\_id} sebagai grup,
agregasi probabilitas ke tingkat anak, bootstrap yang meresampel anak dan bukan citra,
kontrol negatif kebocoran, serta studi degradasi sebagai proxy perangkat sasaran.
\subsection{Instrumen dan produk yang tersedia saat ini}

M-CHAT dan KPSP murah serta mudah disebar, tetapi bergantung pada laporan
manusia~\cite{mardhika2022}. Pada kutub objektif, EarliPoint memperoleh izin FDA 510(k) untuk
anak 16-30 bulan~\cite{fda2022}, tetapi memerlukan eye-tracker khusus di fasilitas klinis
dengan biaya pemeriksaan 599 dolar Amerika~\cite{earlipoint2026}. Pembanding digital terdekat
SenseToKnow: kamera tablet untuk 23 fenotipe perilaku, dengan studi prospektif pada 475
balita di layanan primer AS melaporkan AUC 0,90~\cite{perochon2023}, yang membuktikan
kelayakan kamera perangkat konsumen. Neurogaze mengambil pilihan berbeda: dioperasikan kader
secara luring, memakai stimulus Indonesia, dan menjaga model perilaku tetap terpisah agar
dapat diaudit.

\subsection{Analisis kesenjangan}

Kesenjangan yang kami sasar adalah penerjemahan pengukuran objektif menjadi sistem yang
kelak dapat diuji bersama kader tanpa jaringan. Kamera tablet sudah dipakai penelitian lain,
tetapi validitas Indonesia dan keterpakaian kader belum diuji. Yang kami tambahkan empat hal:
orkestrasi prototipe riset berbasis Posyandu, quality gate yang menahan keputusan ketika
pengukuran gagal, laporan yang dapat diaudit, dan disiplin pelaporan bukti yang menolak
menampilkan angka sebelum datanya ada. Tabel~\ref{tab:beda} pada Lampiran~E merangkumnya.


% =============================================================================
\section{Metodologi}

\subsection{Konsep produk dan alur penggunaan}

Neurogaze menyatu dengan alur SDIDTK yang sudah berjalan, khususnya pada meja keempat
Posyandu. Sasarannya seluruh anak 16-30 bulan yang hadir pada kunjungan terjadwal, dengan
sesi pada usia 18 dan 24 bulan; rentang itu berimpit dengan jendela intervensi paling
produktif sekaligus dengan indikasi instrumen objektif yang sudah diterima
regulator~\cite{fda2022}. Sesi tidak dibatasi pada anak yang positif ceklis, sebab sasarannya
justru anak yang terlewat instrumen berbasis laporan, sehingga waktu per anak dan beban
rujukan masuk sebagai kendala rancangan sejak awal.

Kader mendudukkan anak di depan tablet pada dudukan sederhana, lalu memutar stimulus adegan
sosial yang memancing atensi bergantian antara wajah manusia dan objek non-sosial. Yang
diterima kader adalah kategori tindak lanjut, ringkasan pola pandangan dalam bahasa
sederhana, kualitas pengukuran, dan alasan eksplisit bila sistem menahan keputusan.
Keputusan akhir tetap pada manusia.

\subsection{Arsitektur sistem}

Gambar~\ref{fig:arsitektur} memecah produk menjadi tiga tahap yang berjalan berurutan di
dalam browser, ditambah satu gerbang mutu yang berdiri di antara inferensi dan laporan.

Tahap~A menyiapkan sesi, yaitu pendudukan anak, kalibrasi grid minimal tujuh posisi, lalu
stimulus 96 detik berisi dua belas fase: satu baseline, delapan percobaan isyarat
berskor, satu blok preferential-looking 17 detik untuk Lapis~A, satu blok panggilan nama
13 detik, dan satu penutup. Tahap~B mengubah
video menjadi scanpath lewat rantai pada Lampiran~G, dan di sinilah risiko teknis terbesar
berada; tiap langkahnya punya mode gagal yang berbeda, sehingga quality gate perlu tahu
langkah mana yang gagal untuk memberi kader alasan yang dapat ditindaklanjuti. Tahap~C
menghitung 13 fitur geometri dan menjalankan \texttt{model.json} berisi skaler, koefisien
regresi logistik, dan kalibrator Platt; keluarannya mengisi panel riset yang dijaga penapis
out-of-distribution dan tidak masuk ke jalur keputusan, dengan alasan yang dirinci
Bagian~\ref{sec:lapisan}. Pemeriksaan mutunya sengaja dipecah dua waktu: saat
perekaman aplikasi hanya memeriksa yang dapat diketahui sebelum pemodelan, sedangkan cakupan
fase, sinkronisasi, dan kontrak sesi baru dapat dinilai setelah sesi utuh, sehingga
gerbangnya duduk sesudah inferensi dan berwenang membatalkan keluaran model.

\begin{figure}[t]
\begin{center}
\includegraphics[width=0.82\linewidth]{arsitektur.png}
\end{center}
\caption{Arsitektur produk Neurogaze. Satu sesi 96 detik pada kamera depan tablet melewati
estimasi pandangan, pembentukan scanpath, dan inferensi terkalibrasi, seluruhnya di dalam
browser tanpa jaringan. Quality gate berdiri di antara inferensi dan laporan: sesi yang
gagal memenuhi cakupan fase, deteksi wajah dan iris, sinkronisasi, atau kontrak sesi
menghasilkan status ketiga, yaitu keputusan ditahan disertai alasan yang dapat
ditindaklanjuti kader.}
\label{fig:arsitektur}
\end{figure}

\subsection{Arsitektur keputusan: tiga lapis yang tidak boleh dikaburkan}
\label{sec:lapisan}

Pipeline pada Gambar~\ref{fig:arsitektur} menjelaskan bagaimana angka dihitung, bukan angka
mana yang boleh memutuskan. Keduanya kami pisahkan, sebab mencampurnya adalah cara paling
umum sebuah pengukuran deskriptif berubah menjadi klaim klinis tanpa ada yang memutuskannya.
Sistem ini punya tiga lapis dengan kewenangan berbeda.

\textbf{Lapis~A, GeoPref.} Satu-satunya pemicu rujukan otomatis. Sistem mengukur persentase
waktu pandangan berada pada panel pola geometrik selama blok preferential-looking, lalu
membandingkannya dengan ambang 69 persen dari Wen dkk.~\cite{wen2022}: 1.863 anak usia
12-48 bulan, sensitivitas 17 persen, spesifisitas 98 persen, PPV 81 persen, NPV 65 persen.
Ambangnya tidak kami latih, tidak kami optimasi, dan tidak kami geser. Konsekuensinya harus
dinyatakan sekeras klaimnya: pada sensitivitas 17 persen, alat ini melewatkan sebagian besar
anak autistik, dan pada NPV 65 persen hasil di bawah ambang hampir tidak mengubah keyakinan.
Ini sifat alat \emph{rule-in}, bukan cacat implementasi, dan antarmuka wajib mengatakannya.

\textbf{Lapis~B, profil multi-indeks.} Enam indeks perilaku dilaporkan berdampingan tanpa
skor gabungan: menghadap ke depan, gerak kepala, laju kedip pada fase sosial dan non-sosial,
respons terhadap panggilan nama, dan mengikuti isyarat arah dengan uji tanda dalam sesi.
Keluarga indeksnya mengikuti Perochon dkk.~\cite{perochon2023}, dan tiap indeks membawa AUC
preseden studi itu sebagai provenance di laporan. Menggabungkannya menjadi satu skor
menuntut bobot yang hanya dapat difit pada balita berlabel, yang belum kami punya, sehingga
larangannya ditegakkan tipe: field \texttt{combinedScore} bertipe \texttt{null} dan
\texttt{combinationRuleStatus} adalah literal \texttt{not\_fitted\_requires\_gate\_c}.

\textbf{Lapis~C, model gabungan.} Belum dibangun, karena alasan yang sama.

Regresi logistik Carette berada di luar ketiganya. Ia tetap dikirim bersama aplikasi karena
menjadi dasar seluruh bukti evaluasi pada Bagian~4 dan satu-satunya bukti paritas ujung ke
ujung antara implementasi Python dan TypeScript, tetapi tidak punya jalur kode untuk
memutuskan apa pun; alasannya pada Bagian~\ref{sec:trade-off}.

Tabel~\ref{tab:preseden} menempatkan tiap lapis terhadap preseden yang membenarkannya. Pola
yang kami ikuti sederhana: setiap angka yang berwenang memutuskan harus berasal dari studi
pada populasi dan kelas perangkat yang benar, bukan dari data pengembangan kami sendiri.

\begin{table}[h]
\begin{center}
\small
\begin{tabular}{p{2.1cm}p{3.0cm}p{3.0cm}p{3.4cm}}
\toprule
\textbf{Preseden} & \textbf{Populasi dan perangkat} & \textbf{Performa terbit} &
\textbf{Peran di Neurogaze} \\
\midrule
GeoPref \newline (Wen dkk. 2022) & 1.863 anak, 12-48 bulan; eye-tracker & Sens 17\%,
spec 98\%, PPV 81\%, NPV 65\%, ambang 69\% & Lapis~A. Satu-satunya pemicu rujukan otomatis;
ambang dipakai apa adanya \\
\addlinespace
SenseToKnow \newline (Perochon dkk. 2023) & 475 balita, 17-36 bulan; kamera depan tablet &
AUC 0,90; sens 87,8\%, spec 80,8\% & Lapis~B. Sumber keluarga indeks dan target Gate~C \\
\addlinespace
EarliPoint \newline (FDA 510(k) 2022) & 16-30 bulan; eye-tracker khusus di klinik &
Izin regulator; 599 USD per pemeriksaan & Pembanding kelayakan, bukan sumber ambang \\
\addlinespace
WebGazer \newline (Papoutsaki dkk. 2016; Steffan dkk. 2024) & Webcam; divalidasi
ManyBabies pada 125 anak 18-27 bulan, 16 lab & Galat 4,17$^\circ$; attrition balita 42\% &
Referensi Gate~B dan dasar ekspektasi penyelesaian sesi \\
\bottomrule
\end{tabular}
\end{center}
\caption{Preseden yang menopang tiap lapis keputusan. Tidak satu pun ambang keputusan produk
berasal dari data pengembangan kami: Lapis~A memakai ambang terbit Wen dkk. tanpa penyetelan,
Lapis~B melaporkan indeks secara deskriptif dengan AUC preseden sebagai provenance, dan
target Gate~C diambil dari performa terbit SenseToKnow alih-alih dikarang.}
\label{tab:preseden}
\end{table}

Satu batasan protokol berlaku pada Lapis~A dan harus dibaca bersama ambangnya. Ambang 69
persen sudah diterapkan pada dua tes yang berbeda, dan keduanya tidak berbagi preseden.
Wen dkk.~\cite{wen2022} memvalidasinya pada \emph{GeoPref asli} berdurasi 62,22 detik
(Tobii T120, $1280 \times 1024$, dua AOI $525 \times 363$ piksel) dan melaporkan
sensitivitas 17 persen serta spesifisitas 98 persen pada 1.863 anak usia 12-48 bulan. Moore
dkk.~\cite{moore2018} membawa ambang yang sama ke \emph{Complex Social GeoPref} berdurasi
90 detik---sebagai pilihan konsistensi, bukan hasil optimasi ulang---dan melaporkan
sensitivitas 18 persen, spesifisitas 97 persen, serta AUC 0,74 pada sampel yang jauh lebih
kecil.

Yang berjalan pada aplikasi ini bukan salah satu dari keduanya, melainkan cuplikan 16,75
detik dari video contoh Complex Social yang diterbitkan sebagai Additional file~2 milik
Moore dkk. Cuplikan itu memuat satu dari lima adegan dan berdurasi seperlima protokol
terbitnya. Aplikasi karena itu menyetel \texttt{validatedProtocol} bernilai \texttt{false}
dan \emph{menahan} ambangnya: sesi melaporkan persentase terukur, menyatakan protokolnya
dipersingkat, dan tidak memicu rujukan. Tiap aset stimulus pada
\texttt{app/src/geopref/stimulusMeta.ts} membawa preseden operasionalnya sendiri, sehingga
angka Wen dkk. tidak dapat menempel pada sesi Complex Social maupun sebaliknya.

Tiga sifat aset itu perlu dinyatakan agar tidak terbaca sebagai cacat implementasi.
Pertama, klipnya \emph{tidak bersuara}, dan itu memang protokolnya: metode Moore dkk.
menyatakan bahwa stimulusnya tidak memuat audio. Kedua, panel sosial dan geometriknya
duduk di dalam bingkai hitam $640 \times 360$ dan hanya menempati 19,8 persen luas
bingkai, sedangkan dua AOI pada Wen dkk. menempati 29,1 persen luas layar; presentasi di
perangkat karena itu memangkas kotak hitam agar sudut pandang yang disubtensi mendekati
angka terbit. Ketiga, integritas berkasnya dijaga SHA-256 yang tercatat pada manifes
stimulus.

\subsection{Stimulus blok isyarat: kenapa aktor vektor}
\label{sec:aktor}

Blok isyarat arah memakai aktor yang digambar sebagai SVG dan dianimasikan lewat CSS,
bukan rekaman aktor manusia. Karena blok ini menempati separuh baterai dan memasok satu
dari dua sinyal keputusan yang benar-benar dapat dinilai di lapangan, pilihan itu perlu
dipertahankan dengan bukti.

Yang memberatkan perlu dinyatakan lebih dulu. Tinjauan stimulus eye-tracking untuk joint
attention melaporkan perbedaan antar-kelompok yang lebih tegas pada stimulus beragen
manusia dibandingkan agen animasi. Untuk perbandingan antar-kelompok, agen animasi
kemungkinan memperkecil efek yang hendak diukur. Yang membuat hal itu tidak membatalkan
pilihan ini adalah bentuk sinyalnya: yang kami ambil dari blok ini bukan perbandingan
antar-kelompok, melainkan kontras dalam-subjek---tatapan ke target sesudah isyarat
dibanding sebelum isyarat, pada anak yang sama dan aktor yang sama, di dalam percobaan
yang sama. Uji tanda satu sisi pada delapan percobaan tidak menuntut norma populasi.

Wajah skematik memang memicu \emph{gaze following}. Bayi baru lahir membedakan tatapan
langsung dari tatapan menyamping pada wajah skematik dan berpindah lebih cepat ke target
periferal ketika gerak pupil terlihat~\cite{farroni2004}. Efek itu dibawa oleh sifat
perseptual tertentu, bukan oleh keberadaan wajah secara umum, dan dua di antaranya
merupakan syarat rancangan yang dapat diuji. Pertama, polaritas kontras mata: isyarat
bergantung pada pola lazim pupil gelap di atas sklera terang, dan polaritas terbalik
menghasilkan efek yang jauh lebih lemah~\cite{gredeback2018}. Kedua, gerak pupil harus
terlihat; wajah yang sekadar digambar sudah menghadap ke samping bukan stimulus yang
sama. Aset kami memenuhi keduanya---sklera \texttt{\#fffdf7} melawan pupil
\texttt{\#16181b}, dan bola mata bergeser 16 piksel pada onset dengan transisi 240~ms
berjeda 50~ms sehingga mata berangkat mendahului kepala dan tangan---dan keduanya dijaga
uji kontrak, sehingga penataan ulang gaya visual tidak dapat mencabutnya diam-diam.
Struktur percobaannya sendiri mengikuti paradigma \emph{responding joint attention} pada
balita~\cite{billeci2019}, dengan epok ostensif yang memulihkan kontak mata sebelum ada
informasi arah sama sekali.

Alasan operasionalnya juga bukan sekadar kemudahan. Onset jatuh pada milidetik yang
dideklarasikan protokol di semua tablet tanpa bergantung pada dekoder atau frame yang
jatuh, dan analisis pra/pasca-isyarat bergantung persis pada ketepatan itu. Asetnya kecil
dan berjalan luring penuh, stimulusnya berupa kode berversi yang dapat di-\emph{diff}
alih-alih gumpalan biner, dan tidak ada anak yang perlu direkam sehingga tidak ada
persoalan izin maupun privasi rekaman.

Batasnya tetap harus dinyatakan: tidak ada satu pun bukti bahwa balita merespons aktor
ini. Yang ada adalah paradigma terbit, sifat perseptual yang sesuai literatur, dan
rancangan yang dijaga uji. Apakah balita benar-benar mengikuti isyaratnya adalah
pertanyaan Gate~C.

\subsection{Tidak ada pengukuran yang melintasi batas medium}

Satu sesi menampilkan dua dunia visual: klip video terbit pada blok pilihan tontonan, dan
aktor vektor pada sisanya. Perpindahan itu terlihat oleh peserta, dan pertanyaan yang
tepat bukan apakah tampilannya menyambung melainkan apakah ada besaran yang dihitung
melintasi sambungannya. Preferensi geometrik dihitung seluruhnya di dalam klip oleh
modulnya sendiri; mengikuti isyarat membandingkan pasca-isyarat dengan pra-isyarat di
dalam percobaan vektor yang sama; respons nama adalah deteksi peristiwa di dalam blok
vektor. Tidak satu pun melintasi batas itu.

Satu pengecualian sempat ada dan karena itu dicabut. Diferensial kedipan membandingkan
fase aktor dengan blok video, sehingga kontras ``sosial lawan non-sosial'' tercampur
penuh dengan ``vektor lawan video nyata'', dengan arah bias yang sistematis. Ditambah
jendela 16,75 detik yang mengkuantisasi laju kedip pada kelipatan 3,6 per menit, selisih
yang terukur didominasi derau pencacahan. Indeks itu tetap dilaporkan secara deskriptif
tetapi tidak lagi menjadi sinyal keputusan.


\subsection{Data}

Pengembangan memakai dua dataset publik, satu dipakai dan satu ditolak lewat audit, serta
satu rancangan data prospektif untuk tahap deployment.

\textbf{Dataset scanpath (dipakai).} Autism Eye Tracking Dataset dari Carette
dkk.~\cite{carette2019} berisi 547 citra scanpath beresolusi $640 \times 480$ dari 54
partisipan, terdiri atas 26 anak ASD dan 28 kontrol, berusia rata-rata 7,88 tahun dan direkam
dengan eye-tracker SMI Red-m 250~Hz. Seorang anak menyumbang antara 1 dan 23 rekaman dengan
median 10, sehingga rekaman berulang bukan pengamatan independen dan seluruh evaluasi wajib
dikelompokkan per anak. Dataset ini menjadi anchor validasi algoritmik kami, dan seluruh 54
anaknya dipakai: tiap anak masuk ke lipatan latih, lipatan validasi, dan tepat satu kali ke
lipatan uji tertahan sesuai Bagian~\ref{sec:split}. Ukuran sampelnya kecil untuk standar
epidemiologi, dan kami menjawabnya dengan interval kepercayaan yang meresampel anak.

\textbf{Dataset wajah (ditolak).} Sebuah dataset citra wajah statis publik berisi 2.940 JPEG
berlabel \emph{Autistic} dan \emph{Non\_Autistic}, masing-masing 1.470 citra. Kami menolaknya
meski CNN yang kami latih di atasnya mencapai AUC 0,9324, angka tertinggi di seluruh proyek
ini. Berkasnya bersih secara teknis; yang gagal tata kelolanya, dengan enam dari enam
metadata tidak tersedia. Ketiadaan identitas partisipan saja sudah cukup membuat evaluasi
apa pun tidak dapat ditafsirkan, sebab tidak ada cara membuktikan foto orang yang sama tidak
menyeberang split. Dataset dikarantina dan tidak pernah menyentuh skor Neurogaze. Yang
ditolak dataset ini, bukan modalitas wajah dan kepala. Lampiran~D memuat audit lengkapnya.

\subsection{Pendekatan pemodelan}

Gambar~\ref{fig:pipeline} pada Lampiran~G merangkum enam tahap penelitian, dari citra
scanpath sampai model yang dibekukan, beserta metode yang dipakai di tiap tahap.

Lintasan pandangan dirender sebagai raster $640 \times 480$ lalu diringkas menjadi 13 fitur
geometri tertafsir (Bagian~\ref{sec:tp-fitur}). Enam fitur kinematik Carette tidak dipakai
pada produk: keenamnya bergantung pada kanal warna \emph{velocity-acceleration-jerk} dari
eye-tracker 250~Hz dan tidak dapat direkonstruksi dengan sah dari kamera 30~fps. Sebagai
pembanding representasi, kami juga melatih CNN EfficientNetB0~\cite{tan2019} berpretrain
ImageNet langsung pada citra scanpath, dengan konfigurasi pada Bagian~\ref{sec:tp-model}.

\subsection{Protokol evaluasi: pemisahan latih, validasi, dan uji}
\label{sec:split}

Unit keputusan produk adalah anak, sehingga endpoint utamanya AUC tingkat anak dan seluruh
pemisahan data terjadi pada tingkat anak. Dasarnya terukur: rata-rata 38,2 persen variasi
fitur berada pada tingkat antarpartisipan.

Kami memakai \texttt{GroupKFold} bersarang dengan ID partisipan sebagai grup, memisahkan data
menjadi tiga peran: empat lipatan dalam untuk melatih scaler dan koefisien, satu lipatan
dalam yang ditahan untuk memasang kalibrator Platt dan memilih ambang pada target
sensitivitas 0,9, serta lima lipatan luar sebagai uji tertahan dengan tiap anak muncul tepat
satu kali. Susunan ini membuat setiap angka performa pada Bagian~4 berasal dari anak yang
tidak dilihat model saat dilatih maupun saat dikalibrasi. Assertion otomatis menghentikan
pipeline bila irisan \texttt{participant\_id} antar ketiga peran tidak nol, dan hasilnya nol
di seluruh lipatan; Bagian~\ref{sec:kebocoran} melaporkan apa yang terjadi bila unit
pengelompokan ini salah dipilih.

Ambangnya condong ke sensitivitas tinggi karena asimetri biaya skrining: melewatkan anak
autistik berarti kehilangan jendela intervensi yang kerugiannya sering permanen, sedangkan
merujuk anak tipikal menimbulkan kecemasan dan biaya asesmen yang masih dapat dipulihkan.

\subsection{Implementasi produk}

Neurogaze berjalan sebagai Progressive Web App. MediaPipe Face Landmarker dimuat sebagai WASM
lokal, model statistik disimpan sebagai \texttt{model.json}, dan aset stimulus di-cache
service worker sehingga alur inti berjalan tanpa jaringan setelah kunjungan pertama. Aplikasi
memisahkan empat kontrak: replay deterministik untuk demonstrasi, sesi anak yang melaporkan
pengukuran Lapis~A dan Lapis~B tanpa skor gabungan, Gate~A untuk audit perangkat pada peserta
dewasa, dan halaman bukti Gate~B untuk hasil agreement terhadap WebGazer.js. Laporan sesi
dapat dicetak sebagai serah terima satu halaman untuk Puskesmas, dan log auditnya tinggal di
memori sampai operator memilih mengekspor atau menghapusnya.

\subsection{Lingkup validasi makalah ini dan tahap sesudahnya}
\label{sec:gate}

Produk skrining melewati beberapa jenis validasi, dan mencampuradukkannya adalah sumber klaim
berlebih yang paling umum. Kami memisahkannya jadi empat tingkat: validasi algoritmik menguji
generalisasi ke subjek baru; Gate~A menguji presisi pengukuran perangkat; Gate~B menguji
agreement terhadap WebGazer.js; Gate~C dan D menguji kemampuan skrining dan
operasional pada kohort sasaran. Makalah ini menuntaskan tiga tingkat pertama: validasi
algoritmik, Gate~A, dan Gate~B. Gate~C dan Gate~D tetap terbuka: bukti sekunder dan dasar
prosedur membantu merancang tahap berikutnya, tetapi tidak sama dengan validasi klinis atau
operasional yang selesai. Kriteria seluruh tingkat dicatat pada Lampiran~A agar
sasarannya tidak dapat digeser belakangan.

Kriteria semacam itu mudah ditulis dan mudah dilanggar, jadi kami menegakkannya di kode.
Agregator Gate~B hanya menerima kontrak WebGazer v3, mempertahankan pasangan yang ditahan
dalam denominator, dan menghitung keputusan dari lima kriteria yang tercatat. Halaman bukti
publik \texttt{/validation} membaca snapshot yang angkanya sama dengan ringkasan kanonis,
sedangkan manifest SHA-256 mengunci seluruh 130 ekspor mentah Gate~A dan Gate~B.

% =============================================================================
\section{Hasil}

\subsection{Sinyal geometri scanpath pada tingkat anak}
\label{sec:efek}

Perbedaan antar kelompok muncul dengan ukuran efek besar pada tingkat partisipan
(Gambar~\ref{fig:hasil}a). Delapan dari sebelas fitur dengan efek terbesar berasal dari set
geometri, dipimpin jumlah sel grid aktif ($d = 1{,}69$), fraksi tinta ($d = 1{,}50$), rentang
vertikal ($d = 1{,}48$), dan entropi grid ($d = 1{,}47$). Arahnya konsisten: anak dengan
autisme menyapu bidang layar lebih luas dan lebih merata, sementara anak tipikal memusatkan
pandangan pada zona tengah. Dua fitur kinematik terkuat, kecepatan rata-rata ($d = 1{,}22$)
dan rasio fiksasi ($d = -1{,}15$), berada di bawah empat fitur geometri teratas.

\begin{figure}[t]
\begin{center}
\includegraphics[width=0.80\linewidth]{hasil_utama.png}
\end{center}
\caption{Hasil utama pada 547 scanpath dari 54 anak. (a) Ukuran efek antar kelompok pada
tingkat partisipan; hijau menandai fitur geometri, jingga menandai fitur kinematik yang
bergantung pada encoding eye-tracker, garis titik menandai ambang efek besar $|d| = 0{,}8$.
(b) AUC tingkat anak di bawah proxy degradasi perangkat, dengan pita menunjukkan 95\% CI set
geometri. (c) Jumlah anak yang menyeberang batas split ketika unit pengelompokan salah
dipilih.}
\label{fig:hasil}
\end{figure}

\subsection{Model yang di-deploy: regresi logistik 13 fitur}

Kami membekukan regresi logistik 13 fitur geometri sebagai model MVP, dengan kalibrasi Platt
dipasang pada lipatan validasi dan pengujian pada lipatan luar yang tertahan. Pada data uji
tertahan, model mencapai AUC tingkat anak 0,8228 (95\% CI 0,6974-0,9290) dibandingkan 0,8324
(0,7069-0,9368) untuk set 19 fitur penuh (Tabel~\ref{tab:model} pada Lampiran~E). Selisih berpasangannya
$-0,0137$ dengan 95\% CI $[-0,0593;\ 0,0247]$, jadi keunggulan set penuh pada data bersih
berada di dalam derau. Set geometri menang atau imbang pada empat dari lima keluarga model.
Kalibrasi probabilitasnya wajar, dengan Brier tingkat anak 0,1758 dan kemiringan 1,54.

Target sensitivitas 0,9 melahirkan ambang 0,2243, yang memberi sensitivitas 88,5 persen
dengan spesifisitas 17,9 persen dan rentang bootstrap ambang yang lebar
\mbox{(0,1478-0,4537)}. Spesifisitas serendah itu konsekuensi aritmetika: pada 54 anak,
memaksa sensitivitas mendekati 0,9 menggeser titik kerja ke ujung kurva tempat hampir semua
anak dirujuk. Peringkat model jauh lebih stabil daripada titik kerjanya, sehingga AUC yang
kami perlakukan sebagai hasil dan ambang operasional sebagai parameter per lokasi.

\subsection{Titik kerja, matriks kebingungan, dan analisis kesalahan}
\label{sec:titikkerja}

Satu angka sensitivitas dan satu angka spesifisitas menyembunyikan ruang keputusan yang
sebenarnya tersedia, jadi kami memetakannya. Pada ambang target sensitivitas (0,2243), 46
dari 54 anak dirujuk: model menangkap 23 dari 26 anak ASD tetapi ikut menandai 23 dari 28
kontrol. Pada titik kerja seimbang menurut indeks Youden (0,4985), sensitivitasnya 0,731 dan
spesifisitasnya 0,821 dengan 25 anak dirujuk. Tabel~\ref{tab:ambang} di Lampiran~E merinci
titik-titik di antaranya, termasuk jumlah anak dirujuk per 100 anak diskrining, yaitu besaran
yang menentukan apakah sebuah ambang dapat dijalankan Puskesmas dengan kapasitas rujukan
tertentu. Gambar~\ref{fig:ambang} pada Lampiran~F menampilkan matriks kebingungan kedua titik
kerja itu.

Analisis kesalahan pada titik kerja seimbang berujung pada satu perubahan rancangan produk.
Delapan anak ASD terlewat dan lima kontrol salah ditandai, seluruhnya pada probabilitas
0,169 sampai 0,628, yaitu di sekitar ambang; tidak ada kesalahan berkeyakinan tinggi.
Gambar~\ref{fig:error} pada Lampiran~F menampilkan scanpath-nya, dan polanya terbaca tanpa
statistik: anak ASD yang terlewat punya lintasan pendek dan jarang, sedangkan kontrol yang
salah ditandai punya lintasan yang tidak biasa luas. Angkanya sejalan. Pada 45 anak yang
citranya dapat kami petakan kembali ke rekaman aslinya, korelasi Spearman antara fraksi
tinta dan probabilitas keluaran model adalah 0,927, dan fraksi tinta anak ASD yang terlewat
rata-rata 0,46 persen berbanding 1,94 persen pada yang terdeteksi benar (Mann-Whitney
$p = 0{,}0008$). Keluaran model riset Carette, dengan kata lain, hampir sepenuhnya monoton
terhadap seberapa luas layar disapu pandangan.

Konsekuensinya menyangkut keselamatan. Rekaman yang pendek, terputus, atau tidak terlacak
menghasilkan scanpath jarang, dan scanpath jarang menghasilkan skor rendah, yaitu keluaran
``pantau''. Bila anaknya berisiko, itulah keputusan yang paling berbahaya, dan anak yang
paling sulit direkam bisa jadi anak yang paling mungkin luput. Ambang cakupan pada quality
gate karena itu berfungsi ganda, sebagai penjaga mutu data sekaligus pengaman klinis: sesi
yang tidak memenuhi cakupan minimum wajib ditahan dan diulang, tidak diteruskan menjadi skor
rendah. Lampiran~A menjadikan verifikasinya endpoint wajib pada kohort prospektif.

\subsection{Dampak operasional empat titik kerja}
\label{sec:operasional}

Sensitivitas dan spesifisitas adalah properti kurva; yang menentukan sebuah alat dapat
dijalankan atau tidak adalah berapa banyak anak yang dirujuk dan berapa rujukan yang
dibutuhkan untuk menemukan satu kasus benar. Kami memproyeksikan keempat titik kerja yang
tersedia ke satu antrean Posyandu hipotetis: 1.000 anak, prevalensi 1 persen, dan coverage
teknis 90 persen. Angka coverage-nya asumsi perencanaan yang lebih longgar daripada
attrition 42 persen yang dilaporkan Steffan dkk.~\cite{steffan2024}, dan harus diuji di
lapangan. Tabel~\ref{tab:operasional} memuat hasilnya; artefaknya
\texttt{research/hasil/gate\_c\_simulation.json}, dan lengan regresi logistiknya membaca
\texttt{model.json} sehingga tidak dapat melenceng dari model yang benar-benar dikirim.

\begin{table}[h]
\begin{center}
\small
\begin{tabular}{lrrrrr}
\toprule
\textbf{Titik kerja} & \textbf{Sens} & \textbf{Spec} & \textbf{Rujukan} & \textbf{PPV} &
\textbf{Rujukan/kasus} \\
\midrule
Regresi logistik, sens 0,9      & 0,923 & 0,179 & 82,2\% & 0,011 & 89,1 \\
Regresi logistik, Youden        & 0,731 & 0,821 & 18,4\% & 0,040 & 25,2 \\
\textbf{GeoPref 69\% (aktif)}   & \textbf{0,170} & \textbf{0,980} & \textbf{2,2\%} &
\textbf{0,079} & \textbf{12,6} \\
Target Gate~C (preseden)        & 0,878 & 0,808 & 19,9\% & 0,044 & 22,6 \\
\bottomrule
\end{tabular}
\end{center}
\caption{Proyeksi beban rujukan empat titik kerja pada 1.000 anak, prevalensi 1 persen,
coverage teknis 90 persen. Kolom \emph{Rujukan} adalah persentase anak yang dapat dinilai
dan berakhir dirujuk; kolom \emph{Rujukan/kasus} adalah jumlah rujukan yang dibutuhkan untuk
menemukan satu kasus benar. Simulasi perencanaan, bukan hasil klinis: tidak ada balita
prospektif yang direkrut, dan dua lengan pertama berasal dari 54 partisipan usia sekolah
pada 250~Hz sehingga tidak berlaku untuk balita. Baris ketiga adalah titik kerja yang
benar-benar aktif di produk.}
\label{tab:operasional}
\end{table}

Tabel ini menjelaskan dalam satu baris kenapa alat \emph{rule-in} adalah bentuk yang benar
untuk antrean Posyandu. Titik sensitivitas 0,9 adalah yang paling menggoda dilaporkan, dan ia
merujuk 740 dari 1.000 anak dengan 89 rujukan per satu kasus benar; tidak ada Puskesmas yang
kapasitas rujukannya menyerap itu, dan alat yang menandai empat dari lima anak akan berhenti
dipakai sebelum kasus pertamanya ditemukan. Titik GeoPref merujuk 19 anak pada 12,6 rujukan
per kasus benar. Ia menemukan paling sedikit kasus secara absolut, yaitu 1,5 dari 9 yang ada
pada kohort ini, dan itu harga yang dibayar untuk spesifisitas 98 persen. Pilihan produknya
karena itu bukan pilihan performa melainkan pilihan peran: mendampingi SDIDTK yang menangkap
luas dengan spesifisitas rendah, bukan menggantikannya.

Dua catatan menutup pembacaan tabel ini. Lengan target Gate~C bukan performa Neurogaze
melainkan performa instrumen lain pada kelas perangkat yang sama, ditulis sebagai sasaran
yang harus dikejar. Dan tidak satu pun lengan boleh memakai angka CNN, sebab model itu tidak
dapat dijalankan pada perangkat sasaran, sehingga memasukkannya berarti memproyeksikan biaya
operasional sebuah sistem yang tidak ada.

\subsection{Batas atas representasi: CNN pada citra scanpath}

CNN EfficientNetB0 pada citra scanpath yang sama mencapai AUC tingkat anak 0,8819 (95\% CI
0,7738-0,9681), dengan rata-rata per fold 0,8848 dan simpangan baku 0,1289. Pada ambang 0,476
dari lipatan validasi, balanced accuracy-nya 0,7981 dengan sensitivitas 0,8462 dan
spesifisitas 0,7500, jauh lebih seimbang daripada regresi logistik; 11 dari 54 anak salah
diklasifikasikan, seluruhnya pada probabilitas 0,23 sampai 0,61. Ablasi fine-tuning memberi
hasil nol: membuka 25 persen lapisan teratas hanya mengubah AUC validasi sebesar
$+0,0009 \pm 0,0027$, tidak dapat dibedakan dari derau pada 54 anak. Selisih 0,059 AUC
terhadap model riset Carette adalah harga yang kami bayar dengan sadar; alasannya pada
Bagian~\ref{sec:trade-off}.

\subsection{Verifikasi integritas split}
\label{sec:kebocoran}

Klaim bahwa data uji benar-benar tertahan hanya bermakna bila diuji. Kami menjalankan ulang
pipeline CNN yang identik dan mengganti satu hal saja, yaitu unit pengelompokan split, dari
ID partisipan menjadi klaster duplikat visual berbasis dHash, sebagai kontrol
negatif terhadap protokol Bagian~\ref{sec:split}.

Hasilnya, seperti pada Gambar~\ref{fig:hasil}c, 41 dari 54 anak muncul di sisi latih dan
sisi uji sekaligus, 39 anak menyeberang batas latih/validasi, dan 35 anak menyeberang batas
validasi/uji; dengan pengelompokan per anak ketiganya nol. AUC protokol bocor 0,8351
berbanding 0,8848 pada protokol benar. Arahnya berlawanan dengan dugaan umum, dan pada
simpangan baku antar fold 0,1289 selisih itu tidak dapat dibedakan dari derau, jadi kami
tidak menyimpulkan apa pun tentang besar maupun arah biasnya. Kesimpulan kami lebih
mendasar: ketika 41 anak yang sama muncul di kedua sisi, angkanya berhenti menjawab
pertanyaan yang ingin dijawab produk, yaitu performa pada anak yang belum pernah dilihat.

\subsection{Ketahanan terhadap degradasi perangkat}

Studi degradasi menambahkan derau posisi, sparsifikasi raster sebagai proxy laju sampling,
dropout, dan jitter pose, dengan 200 replikasi bootstrap per kondisi dan tiga seed.
Gambar~\ref{fig:hasil}b dan Tabel~\ref{tab:degradasi} di Lampiran~E memuat hasil lengkapnya.
Terhadap derau, dropout, dan jitter kedua set fitur turun hampir sejajar; terhadap penurunan
laju sampling keduanya berpisah, dan pada kondisi gabungan yang meniru perangkat sasaran
geometri mencapai 0,6827 (95\% CI 0,5384-0,8014) berbanding 0,6049 untuk set penuh.

Hasil itu mengunci 13 fitur geometri untuk MVP. Geometri satu-satunya representasi yang dapat
dihitung dari aliran $(x, y, t)$ 30~fps di perangkat sasaran, dan justru pada kondisi yang
menyerupai perangkat itu ia unggul: selisihnya di sana ($+0,0778$) jauh lebih besar daripada
keunggulan set penuh pada data bersih ($-0,0096$).

\subsection{Hasil sistem dan reproduktibilitas}

Paritas numerik antara implementasi Python dan TypeScript lulus pada toleransi $10^{-12}$,
baik untuk nilai fitur maupun probabilitas model, dengan rangkaian uji hijau penuh. Benchmark
rantai raster sampai inferensi memberi median 4,93~ms pada desktop referensi. Rincian uji,
kontrak luring, dan angka latensi ada pada Lampiran~E.

Analisis desimasi temporal terpisah memakai 27 sesi Gate~B yang membawa stempel waktu nyata.
Saat laju efektif turun dari 26,16 menjadi 13,08~Hz, median drift relatif fitur kinematik
adalah 69,35 persen dan fitur geometri 1,64 persen. Kedua angka ini mengukur drift fitur,
bukan akurasi atau robustness klasifikasi. Enam dari 13 fitur geometri tetap memiliki
Spearman di bawah 0,8 pada 13,08~Hz; pada 6,54~Hz median Spearman geometri turun menjadi
0,4426 dan beberapa fitur menunjukkan drift besar. Hasil ini mendukung pemilihan keluarga
geometri tanpa menyatakan bahwa seluruh fiturnya stabil.

\subsection{Gate~A: validasi instrumen pada kohort dewasa}
\label{sec:gatea}

Gate~A menguji apakah kamera tablet biasa dapat mendeteksi wajah dan iris, membedakan arah
kiri, tengah, dan kanan, menjalankan kalibrasi beserta seluruh fase stimulus, menahan hasil
ketika mutunya gagal, dan tetap bekerja tanpa jaringan. Gerbang ini tidak mengeluarkan skor
risiko dan tidak menguji autisme.

Dua puluh lima peserta dewasa menjalani empat sesi per orang pada tiga jenis perangkat
Android, mencakup kondisi cahaya normal dan redup serta peserta dengan dan tanpa kacamata,
sehingga terkumpul 100 sesi lengkap. Hasilnya 94 dari 100 sesi selesai, galat median
2,2$^\circ$, wajah dan mata terbaca pada 96,4 persen frame, dropout sinyal pandangan 3,6
persen, dan seluruh fungsi inti berjalan luring tanpa satu pun crash. Tabel~\ref{tab:gatea} pada Lampiran~E
merinci sebarannya menurut kondisi pengujian.

Pada 94 dari 100 sesi yang lulus, blok target diketahui menghasilkan konversi sudut lama
dengan median 2,36$^\circ$ dan persentil 90 sebesar 3,58$^\circ$. Jarak pandang tiap sesi
tidak direkam, sehingga angka ini memberi konteks teknis tetapi bukan akurasi absolut eksak.
WebGazer melaporkan 4,17$^\circ$ pada perangkat dan protokol berbeda~\cite{papoutsaki2016};
keduanya bukan uji tanding. Blok target diketahui head-to-head Gate~B belum direkam.

Polanya mengikuti mutu deteksi iris. Cahaya redup dan lensa kacamata sama-sama menurunkan
mutu itu, dan di situlah galat median naik ke kisaran 2,8-2,9$^\circ$ sementara tingkat
keberhasilan sesi turun ke 87-88 persen. Enam sesi tidak menghasilkan laporan sama
sekali, yaitu tiga karena pantulan kacamata, dua karena wajah terlalu miring, dan satu karena
orientasi layar berubah setelah kalibrasi. Seluruh kegagalan itu dikenali sistem dan
diarahkan untuk dikoreksi atau dikalibrasi ulang, bukan diteruskan menjadi skor rendah.

Terhadap kriteria pra-tetap Lampiran~A, keempat ambangnya terpenuhi: galat median
2,2$^\circ$ terhadap batas $\leq 3^\circ$, frame valid 96,4 persen terhadap batas $\geq 85$
persen, penyelesaian sesi 94 persen terhadap batas $\geq 90$ persen, serta nol sesi invalid
yang menghasilkan skor risiko. \textbf{Status Gate~A: lulus.}

\subsection{Gate~B: agreement terhadap WebGazer.js}
\label{sec:gateb}

Gate~B membandingkan aliran gaze Neurogaze dengan implementasi referensi WebGazer.js 3.5.3 yang direkam secara
simultan di browser. Pemilihan pembandingnya bukan soal kemudahan: WebGazer adalah metode
webcam yang divalidasi ManyBabies untuk balita usia 18-27 bulan pada 125 anak di 16
laboratorium~\cite{steffan2024}, sehingga ia satu-satunya referensi terbuka yang punya bukti
pada rentang usia yang mendekati sasaran produk. Studi yang sama melaporkan attrition 42
persen, dan angka itu kami pakai sebagai dasar ekspektasi penyelesaian sesi balita serta
alasan keadaan ditahan dirancang sebagai hasil kelas satu, bukan sebagai kegagalan.
Kontrak referensinya adalah callback
\texttt{setGazeListener} dengan data dan waktu berjalan pada ruang koordinat piksel CSS. Yang
WebGazer bukan ground truth atau standar emas. Yang dibandingkan adalah galat koordinat ternormalisasi, agreement area of interest (AOI), AOI
utama, dan 13 fitur geometri. Pengujian ini mengukur agreement perangkat lunak, bukan
kemampuan mendiagnosis ASD atau kesetaraan dengan eye-tracker laboratorium.

Tiga puluh pasangan terekam. Dua puluh tujuh berstatus siap dan tiga ditahan, sehingga valid
pair rate 90 persen dan retry rate 10 persen. Galat median antarliran adalah 44,159 piksel
atau 0,040997 pada skala ternormalisasi. Seluruh jarak ini direderivasi dari koordinat mentah
oleh \texttt{research/recompute\_gate\_b.py} dan tereproduksi sampai 0,001 piksel. Agreement AOI
rata-rata hasil rekomputasi adalah 99,7118 persen; nilai yang ditulis browser saat perekaman
99,7574 persen, dan selisihnya berasal dari empat pasangan yang masing-masing memiliki satu
sampel tepat di luar kotak AOI menurut \texttt{neurogaze-aoi-v3.1.0}. Selisih itu diterbitkan di
\texttt{gate\_b\_summary.json} alih-alih didamaikan. AOI utama cocok pada seluruh 27 pasangan siap.
Angka agreement setinggi ini perlu dibaca dengan hati-hati: kotak AOI selebar 28 persen layar
sementara galat median antaraliran hanya 4,1 persen lebar layar, sehingga kesepakatan tinggi
nyaris tidak terhindarkan dan bukan bukti akurasi. Rata-rata ICC(A,1) untuk 13 fitur adalah
0,505043; nilai ini dilaporkan secara deskriptif dan bukan kriteria kelulusan.

Seluruh lima kriteria yang tercatat terpenuhi: 30 pasangan terhadap batas $\geq 30$, valid
pair rate 90 persen terhadap batas $\geq 90$ persen, galat median ternormalisasi 0,040997
terhadap batas $\leq 0,05$, agreement AOI rata-rata 99,7118 persen terhadap batas $\geq 95$
persen, dan agreement AOI utama 100 persen terhadap batas $\geq 95$ persen.
\textbf{Status Gate~B: lulus.}

Kesimpulan ini terbatas pada agreement terhadap implementasi referensi WebGazer.js dalam kohort yang terekam.
Generalisasi klinis pada balita, sensitivitas, dan spesifisitas tetap harus dinilai melalui
Gate~C prospektif dengan hasil klinis independen.

\subsection{Kontrol positif dewasa dan kesiapan}

Kontrol positif menguji apakah rantai kamera-ke-angka merespons ketika orang dewasa sengaja
memproduksi pola yang dicari. Dua belas orang dewasa memberi persetujuan dan menghasilkan
23 sesi pada tiga perangkat; 15 sesi lulus mutu. Denominator per kondisi adalah 9 dari 11
sesi menonton biasa dan 6 dari 12 sesi pola diproduksi yang dapat dipakai. Dalam mode
demonstrasi, aturan menyala pada 0 dari 9 sesi biasa dan 4 dari 6 sesi pola diproduksi.
Mode ini mengunci \texttt{emitsReferral} pada \texttt{false}. Tiga sinyal terpisah tanpa
tumpang tindih, tetapi AUC 1,00 hanya menyatakan tidak ada pasangan yang terurut salah;
margin terdekat dan denominator lengkap lebih informatif. Hasil ini adalah manipulation
check dan responsivitas end-to-end, bukan sensitivitas, spesifisitas, ASD, atau status
klinis peserta.

Matriks kesiapan yang dihasilkan bersama artefak publik menempatkan rantai pengukuran di
perangkat sebagai siap untuk demonstrasi rekayasa dan respons dewasa sebagai telah
ditunjukkan. Rujukan otomatis balita ditahan. Keterpakaian kader dan validitas balita
Indonesia belum diuji. Validasi prospektif memerlukan mitra yang mampu menjalankan kaji
etik, izin orang tua, acuan klinis buta, linkage privat, rekrutmen, analisis fairness dan
kegagalan, serta validasi prospektif sebelum pemilihan titik operasi.

\subsection{Ringkasan capaian validasi}

Tabel~\ref{tab:gate} pada Lampiran~E merangkum posisi tiap tingkat validasi terhadap lingkup
makalah ini. Tiga tingkat sudah tuntas. Validasi algoritmiknya menjawab seluruh endpoint yang
kami pra-tetapkan dari data uji tertahan. Gate~A lulus pada 100 sesi dewasa lintas tiga
perangkat, dua kondisi pencahayaan, serta peserta dengan dan tanpa kacamata. Gate~B lulus
pada 30 perbandingan browser simultan terhadap WebGazer.js, dengan seluruh lima kriteria
agreement terpenuhi.

Dua tingkat sisanya tetap terbuka. Untuk Gate~C, targetnya kami ambil dari literatur alih-alih
dari data kami sendiri: sensitivitas 87,8 persen dan spesifisitas 80,8 persen yang dicapai
SenseToKnow pada kamera depan tablet untuk anak 17-36 bulan~\cite{perochon2023}. Angka
sekunder dari dataset Carette sengaja tidak diwariskan sebagai kandidat titik kerja, sebab
populasi maupun instrumennya berbeda dan CNN-nya tidak dapat dijalankan pada perangkat
sasaran. Cabang prospektifnya tetap terbuka dengan kohort balita nol. Untuk Gate~D, protokol ramah anak, jarak layar
sekitar 60~cm, posisi duduk fleksibel, kalibrasi lima titik, dan stimulus naturalistik sudah
dipakai pada studi eye-tracking anak, tetapi pengujian lapangan tetap terbuka dengan nol
Posyandu, nol kader, dan nol sesi anak. Bukti sekunder ini tidak setara dengan validasi
klinis atau operasional Neurogaze.

% =============================================================================
\section{Nilai Bisnis dan Jalur Adopsi}

\subsection{Peran, biaya, dan titik masuk layanan}

Tiga peran sering tertukar, jadi kami pisahkan sebagai hipotesis implementasi. Penerima
manfaat yang dituju adalah anak 16--30 bulan beserta pengasuhnya. Operator yang diusulkan
adalah kader di bawah supervisi Puskesmas, tetapi belum ada kader yang menguji aplikasi.
Skema pembiayaan dan penerimaan laporan juga belum diuji dengan Dinas Kesehatan atau
praktisi.

PWA statis membuat biaya distribusi perangkat lunak mendekati nol dan tidak memerlukan
lisensi per kursi. Biaya operasi belum diukur: waktu kader, pelatihan, dukungan, pencetakan,
dudukan, pengulangan, pemeliharaan, penggantian tablet, dan tindak lanjut klinis semuanya
tetap ada. Karena itu simulasi biaya per sesi maupun per kasus di Lampiran hanya skenario
bersyarat apabila protokol penuh dan titik operasi terbit berhasil direplikasi, bukan
performa atau biaya NeuroGaze yang diamati. Sebagai pembanding, EarliPoint, instrumen tatapan
berizin FDA untuk rentang usia yang sama, memerlukan eye-tracker khusus di fasilitas klinis
dengan biaya pemeriksaan 599 dolar Amerika~\cite{earlipoint2026}. Neurogaze tidak bersaing
dengan ketepatan instrumen itu; pesaingnya ketiadaan pemeriksaan sama sekali.

Titik masuk yang diusulkan adalah alur Posyandu. Baterai pengukuran berlangsung 67 detik;
izin, penempatan, kalibrasi, kemungkinan pengulangan, dan penjelasan menambah waktu
kunjungan. Durasi total dan penerimaannya belum diukur pada kader atau balita.

\subsection{Skalabilitas dan jalur adopsi bertahap}

Skala kabupaten belum dapat diproyeksikan dari data proyek ini. Kapasitas asesmen lanjutan,
attrition, waktu kader, dan pengulangan harus diukur sebelum jumlah sesi atau biaya dapat
dipakai untuk keputusan program.

Perluasannya bertahap, dengan gerbang validasi sebagai syarat lolos tiap tahap, bukan target
komersial. Tahap pertama kemitraan dengan rumah sakit, klinik tumbuh kembang, dan pusat
terapi, yang menyelesaikan dua masalah sekaligus: mitra klinis menyediakan kohort berlabel,
klinisi yang dibutakan untuk menetapkan acuan, serta tata kelola kaji etik untuk Gate~C.
Kemitraan ini tidak dapat direduksi menjadi persetujuan administratif. Tahap kedua
adopsi oleh Dinas Kesehatan kabupaten setelah Gate~D menunjukkan produk dapat dijalankan
kader tanpa mengganggu layanan. Tahap ketiga perluasan lintas wilayah dengan kalibrasi ulang
terhadap norma setempat, karena norma pola tatapan tidak otomatis berpindah antar populasi.

% =============================================================================
\section{Diskusi}

\subsection{Cakupan kesimpulan yang ditopang hasil}

Validasi yang kami tuntaskan menopang lima kesimpulan. Geometri scanpath memuat sinyal
berukuran efek besar yang bertahan pada anak yang tidak pernah dilihat model (AUC 0,8228),
dan bertahan lebih baik daripada fitur kinematik pada kondisi yang menyerupai kamera
konsumen. Protokol splitnya terverifikasi tahan bocor lewat kontrol negatif. Seluruh
rantai dari kalibrasi sampai laporan berjalan luring dengan paritas numerik $10^{-12}$
terhadap implementasi riset, sehingga angka di notebook dan produk yang dipakai kader dijamin
identik oleh uji otomatis. Kamera tablet biasa cukup presisi untuk mengukur arah pandangan
pada kondisi lapangan yang wajar, dengan galat median 2,2$^\circ$ dan penyelesaian sesi 94
persen pada 100 sesi (Bagian~\ref{sec:gatea}). Aliran gaze Neurogaze juga memenuhi seluruh
kriteria agreement terhadap WebGazer.js, dengan galat median ternormalisasi 0,040997 dan
agreement AOI 99,7118 persen pada 27 dari 30 pasangan yang siap (Bagian~\ref{sec:gateb}).

Batas kelimanya sama dan perlu dinyatakan tegas: Gate~A membuktikan kelayakan teknis dan
Gate~B membuktikan agreement terhadap WebGazer.js, tetapi tidak satu pun membuktikan akurasi diagnosis
ASD. Klaim klinis tetap menunggu kohort balita dengan hasil klinis independen yang dinilai
secara buta.

\subsection{Alasan model ber-AUC tertinggi tidak dipakai di produk}
\label{sec:trade-off}

CNN scanpath unggul 0,059 AUC dan titik kerjanya jauh lebih seimbang, tetapi tidak masuk
produk. Penghalang pertamanya kontrak input. Model yang di-deploy harus memakan aliran
$(x, y, t)$ dari kamera 30~fps di dalam browser, sedangkan citra scanpath Carette adalah
raster yang sudah dirender dengan kanal warna kecepatan, akselerasi, dan jerk dari sinyal
250~Hz. Melatih CNN pada citra itu lalu memberinya raster dari kamera tablet berarti menguji
model pada distribusi yang berbeda dari distribusi pelatihannya, dan kenaikan 0,059 pada
dataset sumber tidak menjamin apa pun setelah perpindahan itu.

Penghalang kedua bersifat operasional. Regresi logistik 13 fitur berjalan pada median
4,93~ms, dapat diekspor sebagai \texttt{model.json} berisi koefisien dan kalibrator yang
terbaca manusia, dan kontribusi tiap fiturnya dapat diaudit. Yang menjalankan produk ini
kader, dengan latar pelatihan berbeda dari klinisi, dan di konteks itu kemampuan menjelaskan
alasan seorang anak disarankan dirujuk bernilai lebih besar daripada 0,059 AUC pada dataset
yang populasinya berbeda dari sasaran. Angka CNN tetap kami laporkan karena menjawab
pertanyaan lain yang sah, yaitu berapa banyak sinyal yang ada di representasi ini bila
kapasitas model tidak dibatasi.

\subsection{Menolak dataset yang menguntungkan}
\label{sec:tolak-wajah}

CNN yang kami latih pada dataset wajah statis sebagai audit metodologis mencapai AUC uji
0,9324, angka tertinggi di proyek ini. Kami tidak memasukkannya sebagai hasil, dengan tiga
alasan yang berdiri sendiri-sendiri. Inferensi karakteristik neurologis dari morfologi wajah
statis memikul beban sejarah fisiognomi. Dataset itu tidak punya persetujuan subjek
terdokumentasi untuk anak-anak yang dapat diidentifikasi. Dan tanpa identitas partisipan,
tidak ada evaluasi yang dapat membuktikan model tidak sekadar menghafal orang; alasan ketiga
ini teknis, jadi menambah izin sekalipun tidak akan memperbaikinya.

Uji shortcut menunjukkan performa itu tidak seluruhnya artefak teknis, tetapi juga tidak
bersih: properti berkas saja mencapai AUC 0,6757 dan penambahan statistik piksel tingkat
rendah menaikkannya ke 0,7515. Sisa jaraknya tidak terjelaskan, dan tanpa provenance kami
tidak punya cara mengetahui asalnya. Keputusan menolak diambil di tingkat tata kelola, bukan
tingkat metrik; Lampiran~D memuat audit lengkapnya.

\subsection{Prasyarat kelembagaan untuk tahap berikutnya}
\label{sec:kendala}

Validasi algoritmik, Gate~A, dan Gate~B dapat kami tuntaskan sendiri karena ketiganya tidak
menuntut peserta anak dengan diagnosis. Gate~C prospektif dan Gate~D lapangan berbeda:
keduanya terhalang prasyarat kelembagaan di luar jangkauan teknis kami, dan kami mengujinya
langsung. Untuk memperoleh data sesi dari anak dengan diagnosis, kami menghubungi lima
lembaga terapi dan layanan autisme di Surabaya dan sekitarnya. Tidak satu pun membuka akses,
dan itu keputusan yang benar dari pihak mereka: lembaga terapi memikul kewajiban perlindungan
terhadap anak dan keluarga, dan permintaan pengumpulan data dari pihak tanpa tinjauan etik
formal, sponsor institusional, maupun rekam jejak klinis memang seharusnya ditolak.

Jalur data klinis, dengan kata lain, hanya terbuka lewat kemitraan institusional berizin
etik, jadi prioritas pertama pada Bagian~\ref{sec:lanjut} kemitraan klinis, bukan pekerjaan
teknis tambahan.

\subsection{Risiko dan mitigasi}
\label{sec:risiko}

Risiko pelabelan ditekan dengan membatasi keluaran pada rekomendasi rujukan, menampilkan
ketidakpastian, dan mensyaratkan persetujuan yang dapat ditarik. Risiko privasi ditekan
dengan memproses seluruh video di perangkat dan hanya mempertahankan fitur serta skor
pseudonim. Risiko klaim berlebih ditekan lewat struktur aplikasinya: sesi kamera melaporkan pengukuran
Lapis~A dan Lapis~B beserta provenance tiap indeks dan tidak pernah mengeluarkan skor
gabungan maupun probabilitas ASD, keluaran regresi logistik dikurung di panel riset berpenapis
out-of-distribution, dan halaman bukti publik menolak menampilkan metrik tanpa berkas
comparison dan ambang pra-registrasi. Lapis~A adalah satu-satunya kandidat rujukan dengan
titik operasi terbit, tetapi klip 16,75 detik menahan keluarannya. Tabel~\ref{tab:operasional}
hanya mensimulasikan beban jika titik operasi protokol penuh kelak direplikasi.

Dua risiko masih terbuka. Performa per subkelompok memerlukan demografi yang tidak disediakan
dataset anchor, dan perilaku model terhadap anak dengan hambatan perkembangan non-autisme
hanya terukur lewat kontrol negatif sulit dari kohort prospektif. Keduanya masuk sebagai
endpoint wajib pada Lampiran~A.

\subsection{Catatan tata kelola: keadilan, pra-registrasi subgrup, dan konteks Indonesia}
\label{sec:tatakelola}

Tiga hal perlu dinyatakan sebelum sistem ini dibawa ke lapangan, dan kami menyatakannya di
muka karena ketiganya lebih mudah diabaikan daripada dibantah.

\emph{Pertama, preseden yang kami pakai punya kesenjangan performa lintas kelompok.} Pada
studi SenseToKnow, spesifisitasnya 53,6 persen pada anak kulit hitam berbanding 82,7 persen
pada anak kulit putih~\cite{perochon2023}. Kesenjangan sebesar itu pada instrumen yang
menjadi preseden Lapis~B berarti risiko rujukan berlebih tidak terdistribusi rata, dan pada
konteks Indonesia sumber variasinya bisa berbeda lagi: warna kulit, bentuk kelopak, dan
pemakaian kerudung yang menutupi tepi wajah semuanya menyentuh langkah deteksi landmark.
GeoPref sendiri dilaporkan setara lintas ras pada studi aslinya, tetapi kesetaraan pada satu
populasi bukan jaminan pada populasi lain.

\emph{Kedua, analisis subgrup harus dipra-registrasi, bukan dicari setelah datanya terlihat.}
Kami memasukkan performa per subkelompok sebagai endpoint wajib Gate~C pada Lampiran~A,
dengan stratifikasi yang ditetapkan lebih dulu, sebab subgrup yang dipilih setelah melihat
hasil selalu dapat dibuat tampak baik. Ini juga satu-satunya cara menutup batasan yang
diwarisi dari data pengembangan: dataset anchor tidak menyediakan demografi sama sekali,
sehingga tidak ada analisis keadilan yang dapat dilakukan pada hasil di Bagian~4.

\emph{Ketiga, tidak satu pun instrumen yang dipakai di sini divalidasi di Indonesia.}
Ambang GeoPref berasal dari kohort Amerika Serikat, keluarga indeks Lapis~B dari layanan
primer Amerika Serikat, dan referensi Gate~B dari 16 laboratorium yang seluruhnya di luar
Indonesia. Yang berpindah lintas populasi bukan hanya distribusi perilakunya, tetapi juga
jalur layanan sesudah rujukan: sebuah alat \emph{rule-in} hanya bernilai bila ada kapasitas
asesmen yang menerima rujukannya. Karena itu Gate~D mengukur beban layanan, bukan hanya
kegunaan antarmuka. Dan perlu diulang sekali lagi di bagian ini, karena inilah batasan yang
paling mudah lupa disebut: belum ada satu pun balita dalam bukti mana pun di makalah ini.

\subsection{Pengembangan lanjutan}
\label{sec:lanjut}

Dengan Gate~A dan Gate~B tertutup, sisa pekerjaannya tidak lagi bersifat teknis melainkan
kelembagaan, dan urutan di bawah mengikuti prasyarat itu. Pertama, kemitraan formal dengan
rumah sakit atau klinik tumbuh kembang berkomite etik, satu-satunya kunci yang membuka Gate~C
prospektif. Kedua, kohort prospektif Gate~C pada balita 16-30 bulan dengan acuan klinis
buta dan kontrol negatif sulit, yaitu anak dengan hambatan perkembangan non-autisme; hanya di
tahap inilah ambang klinis boleh dipilih. Sasarannya sensitivitas 87,8 persen dan
spesifisitas 80,8 persen, yaitu performa terbit SenseToKnow pada kelas perangkat yang
sama~\cite{perochon2023}, dan endpoint keadilan pada Bagian~\ref{sec:tatakelola} masuk ke
protokolnya sejak awal. Ketiga, Gate~D di Posyandu dengan target 5 Posyandu, 20 kader, 200
sesi anak, dan tiga tipe tablet Android, untuk mengukur waktu layanan, keberhasilan sesi
pertama, pemahaman laporan, dan penerimaan orang tua. Keempat, perluasan matriks Gate~A dan
Gate~B ke lebih banyak tipe perangkat dan kondisi klinis, karena kedua gerbang itu lulus pada
sampel yang masih kecil. Cabang vokalisasi dan gerak kepala diuji bersama kohort prospektif
sebagai hipotesis.

% =============================================================================
\section{Kesimpulan}

Neurogaze saat ini adalah demonstrasi rekayasa dengan rantai pengukuran, penolakan, dan
pelaporan yang dapat diaudit. Model Carette memberi bukti konsep participant-grouped pada
domain sumber anak usia sekolah Prancis 250~Hz; ia tidak memiliki uji eksternal, sengaja
ditolak penjaga OOD pada sesi sekarang, dan tidak pernah menentukan keputusan. Gate~A
mencakup 100 sesi berulang dari 25 dewasa dengan 94 lulus mutu. Gate~B mencakup 30
perbandingan simultan terhadap implementasi referensi WebGazer.js, dengan 27 siap dan tiga
ditahan; ia mengukur agreement perangkat lunak, bukan ground truth. Kontrol positif dewasa
menunjukkan respons kamera-ke-angka pada 15 dari 23 sesi yang lulus mutu, dengan aturan demo
0/9 pada kondisi biasa dan 4/6 pada pola diproduksi. Rujukan otomatis balita ditahan karena
klip 16,75 detik tidak mereplikasi protokol penuh titik operasi 69 persen.

Belum ada balita, uji kader, validasi Indonesia, wawancara praktisi, atau linkage klinis.
Tahap berikutnya memerlukan mitra yang mampu menjalankan kaji etik, izin orang tua, acuan
klinis buta, linkage privat, rekrutmen, analisis fairness dan kegagalan, serta validasi
prospektif sebelum titik operasi dipilih. Tidak satu pun hasil makalah ini merupakan bukti
akurasi skrining ASD pada NeuroGaze.

\subsubsection*{Ketersediaan artefak}

\begin{sloppypar}
Kode dan notebook ada di
\url{https://github.com/raphaeladikara/Dataton_Project_Rizzz}, dan aplikasinya dapat dicoba
tanpa pemasangan di \url{https://dataton-project-rizzz.vercel.app}. Bobot model gaze
yang terlatih tersedia di
\url{https://huggingface.co/raphaeladikara/neurogaze-gaze-lr-v1}.
\end{sloppypar}

% =============================================================================
\newpage
\small
\begin{thebibliography}{99}

\bibitem{dawson2010}
G.~Dawson, S.~Rogers, J.~Munson, M.~Smith, J.~Winter, J.~Greenson, A.~Donaldson, dan
J.~Varley. Randomized controlled trial of an intervention for toddlers with autism: the Early
Start Denver Model. \emph{Pediatrics}, 125(1):e17-e23, 2010.

\bibitem{vanthof2021}
M.~van~'t Hof, C.~Tisseur, I.~van Berckelaer-Onnes, A.~van Nieuwenhuyzen, A.~M. Daniels,
M.~Deen, H.~W. Hoek, dan W.~A. Ester. Age at autism spectrum disorder diagnosis: a systematic
review and meta-analysis from 2012 to 2019. \emph{Autism}, 25(4):862-873, 2021.

\bibitem{harper2026}
A.~L. Harper, D.~E. Sarver, dan K.~F. Sly. Validation of the Modified Checklist for Autism in
Toddlers (M-CHAT): a replication study of diagnostic accuracy. \emph{Journal of Autism and
Developmental Disorders}, 2026. doi:10.1007/s10803-026-07330-3.

\bibitem{sturner2016}
R.~Sturner, B.~Howard, P.~Bergmann, D.~Morrel, L.~Andon, C.~Marks, P.~Rao, dan M.~Landa.
Autism screening with online decision support by primary care pediatricians aided by
M-CHAT/F. \emph{Pediatrics}, 138(3):e20153036, 2016.

\bibitem{fathmawati2023}
Fathmawati, L.~Andari, J.~Saimin, D.~F. Damayanti, Z.~Akhmadi, dan E.~D. Septiawan.
Distribution of autistic children in Pontianak, West Kalimantan Province, Indonesia.
\emph{Malaysian Journal of Public Health Medicine}, 23(2):124-130, 2023.
doi:10.37268/mjphm/vol.23/no.2/art.1826.

\bibitem{aditya2021}
C.~J. Aditya, J.~K. Dahliana, A.~D. Widodo, dan R.~Sekartini. Autism spectrum disorder
screening in children aged 16-30 months using the Modified Checklist for Autism in
Toddlers-Revised (M-CHAT-R). \emph{Paediatrica Indonesiana}, 61(5):247-252, 2021.
doi:10.14238/pi61.5.2021.247-52.

\bibitem{kemenkopmk2023}
Kementerian Koordinator Bidang Pembangunan Manusia dan Kebudayaan RI. Posyandu miliki peran
penting ciptakan generasi emas Indonesia, 2023.

\bibitem{mardhika2022}
A.~Mardhika, J.~Susanto, E.~Qurniyawati, dan A.~P.~M. Tyas. Empowerment of healthcare cadres
on stimulation of early detection and intervention of growth and development. \emph{Abdimas
Unwahas}, 7(2), 2022. doi:10.26905/abdimas.v7i2.6909.

\bibitem{jones2013}
W.~Jones dan A.~Klin. Attention to eyes is present but in decline in 2-6-month-old infants
later diagnosed with autism. \emph{Nature}, 504(7480):427-431, 2013.

\bibitem{pierce2011}
K.~Pierce, D.~Conant, R.~Hazin, R.~Stoner, dan J.~Desmond. Preference for geometric patterns
early in life as a risk factor for autism. \emph{Archives of General Psychiatry},
68(1):101-109, 2011.

\bibitem{moore2018}
A.~Moore, M.~Wozniak, A.~Yousef, C.~C. Barnes, D.~Cha, E.~Courchesne, dan K.~Pierce.
The geometric preference subtype in ASD: identifying a consistent, early-emerging
phenomenon through eye tracking. \emph{Molecular Autism}, 9:19, 2018.
doi:10.1186/s13229-018-0202-z.

\bibitem{wen2022}
T.~H. Wen, A.~Cheng, C.~Andreason, J.~Zahiri, Y.~Xiao, R.~Xu, B.~Bao, E.~Courchesne,
C.~C. Barnes, S.~J. Arias, dan K.~Pierce. Large scale validation of an early-age
eye-tracking biomarker of an autism spectrum disorder subtype. \emph{Scientific Reports},
12(1):4253, 2022. doi:10.1038/s41598-022-08102-6.

\bibitem{billeci2019}
L.~Billeci, S.~Narzisi, A.~Tonacci, dkk. Disentangling the initiation from the response
in joint attention: an eye-tracking study in toddlers with autism spectrum disorders.
\emph{Frontiers in Psychology}, 10:2187, 2019. doi:10.3389/fpsyg.2019.02187.

\bibitem{farroni2004}
T.~Farroni, S.~Massaccesi, D.~Pividori, dan M.~H. Johnson. Gaze following in newborns.
\emph{Infancy}, 5(1):39-60, 2004.

\bibitem{gredeback2018}
G.~Gredeb\"ack dkk. Schematic eye-gaze cues influence infants' object encoding dependent
on their contrast polarity. \emph{Scientific Reports}, 7:9506, 2017.
doi:10.1038/s41598-017-10108-4.


\bibitem{fda2022}
U.S. Food and Drug Administration. EarliPoint System, 510(k) K213882, 2022; EarliPoint,
510(k) K230337, 2023.

\bibitem{earlipoint2026}
EarliPoint Health. Frequently asked questions. Diakses 2026.

\bibitem{perochon2023}
S.~Perochon, J.~M. Di~Martino, K.~L.~H. Carpenter, S.~Compton, N.~Davis, B.~Eichner,
S.~Espinosa, L.~Franz, P.~R. Krishnappa~Babu, G.~Sapiro, dan G.~Dawson. Early detection of
autism using digital behavioral phenotyping. \emph{Nature Medicine}, 29:2489-2497, 2023.
doi:10.1038/s41591-023-02574-3.

\bibitem{steffan2024}
A.~Steffan, L.~Zimmer, N.~Arias-Trejo, M.~Bohn, R.~Dal~Ben, dkk. Validation of an
open source, remote web-based eye-tracking method (WebGazer) for research in early
childhood. \emph{Infancy}, 29(1):31-55, 2024. doi:10.1111/infa.12564.

\bibitem{papoutsaki2016}
A.~Papoutsaki, P.~Sangkloy, J.~Laskey, N.~Daskalova, J.~Huang, dan J.~Hays. WebGazer:
scalable webcam eye tracking using user interactions. Dalam \emph{Proceedings of the
25th International Joint Conference on Artificial Intelligence (IJCAI)}, halaman
3839-3845, 2016.

\bibitem{carette2019}
R.~Carette, M.~Elbattah, F.~Cilia, G.~Dequen, J.-L. Gu\'erin, dan J.~Bosche. Learning to
predict autism spectrum disorder based on the visual patterns of eye-tracking scanpaths.
Dalam \emph{Proceedings of BIOSTEC 2019 (HEALTHINF)}, halaman 103-112, 2019.

\bibitem{tan2019}
M.~Tan dan Q.~V. Le. EfficientNet: rethinking model scaling for convolutional neural
networks. Dalam \emph{Proceedings of the 36th International Conference on Machine Learning},
halaman 6105-6114, 2019.

\end{thebibliography}
\normalsize

% =============================================================================
\newpage
\appendix
\section*{Lampiran A. Kriteria validasi pra-tetap}
\addcontentsline{toc}{section}{Lampiran A}

Bagian pertama tabel memuat endpoint validasi algoritmik yang menjadi lingkup makalah ini;
seluruhnya ditetapkan sebelum eksperimen dijalankan dan seluruhnya terjawab dari data uji
tertahan. Bagian kedua memuat kriteria Gate~A dan Gate~B, yang juga ditetapkan lebih dulu dan
kini seluruhnya terjawab dari data pengujian instrumen. Bagian ketiga memuat kriteria Gate~C
dan Gate~D yang belum terjawab, dipublikasikan lebih dulu justru
agar sasarannya tidak dapat digeser setelah datanya terlihat. Kegagalan pada satu baris
menghentikan konfigurasi perangkat atau stimulus yang diuji; nilai di zona tengah mewajibkan
satu putaran perbaikan terarah lalu pengukuran ulang pada sampel baru, bukan penyetelan
ulang model terhadap data yang sama.

\begin{table}[h]
\begin{center}
\small
\begin{tabular}{llll}
\toprule
\textbf{Tingkat} & \textbf{Endpoint} & \textbf{Ambang lulus} & \textbf{Hasil} \\
\midrule
\multicolumn{4}{l}{\emph{Validasi algoritmik: lingkup makalah ini}} \\
Algoritmik & AUC tingkat anak, uji tertahan & $> 0,75$ & 0,8228 \checkmark \\
Algoritmik & irisan anak antar split         & $= 0$    & 0 \checkmark \\
Algoritmik & AUC pada proxy tablet gabungan  & $> 0,60$ & 0,6827 \checkmark \\
Algoritmik & set fitur dapat dihitung 30~fps & wajib    & 13 fitur \checkmark \\
Algoritmik & paritas Python-TypeScript      & $< 10^{-9}$ & $10^{-12}$ \checkmark \\
Algoritmik & tata kelola dataset terdokumentasi & wajib & audit lengkap \checkmark \\
\midrule
\multicolumn{4}{l}{\emph{Gate~A dan B: dipra-registrasi, sudah terjawab}} \\
A & galat median kalibrasi        & $\leq 3^\circ$        & 2,2$^\circ$ \checkmark \\
A & frame wajah/iris valid        & $\geq 85\%$           & 96,4\% \checkmark \\
A & session completion            & $\geq 90\%$           & 94\% \checkmark \\
A & skor dari sesi invalid        & $= 0$                 & 0 \checkmark \\
A & P90 ekstraksi                 & $< 100$~ms            & 7,67~ms \checkmark \\
B & jumlah pasangan               & $\geq 30$             & 30 \checkmark \\
B & valid pair rate               & $\geq 90\%$           & 90\% \checkmark \\
B & galat median ternormalisasi   & $\leq 0,05$           & 0,040997 \checkmark \\
B & AOI agreement rata-rata       & $\geq 95\%$           & 99,7118\% \checkmark \\
B & agreement AOI utama           & $\geq 95\%$           & 100\% \checkmark \\
\midrule
\multicolumn{4}{l}{\emph{Gate~C dan D: dipra-registrasi, belum terjawab}} \\
C & AUC tingkat anak balita       & ditetapkan bersama klinisi & kohort prospektif \\
C & kontrol negatif sulit         & wajib ada             & kohort prospektif \\
C & sesi cakupan rendah ditahan   & wajib                 & kohort prospektif \\
D & median waktu layanan          & $\leq 5$ menit        & studi lapangan \\
D & kapasitas per perangkat       & $\geq 10$ anak/jam    & studi lapangan \\
\bottomrule
\end{tabular}
\end{center}
\end{table}

\section*{Lampiran B. Tiga status laporan produk}
\addcontentsline{toc}{section}{Lampiran B}

\begin{figure}[h]
\begin{center}
\includegraphics[width=\linewidth]{produk.png}
\end{center}
\caption{Tiga status laporan yang dapat diuji melalui replay deterministik bawaan.
(a) rekomendasi rujukan; (b) pantau sesuai alur SDIDTK; (c) keputusan ditahan karena
kualitas pengukuran tidak memadai, disertai alasan yang dapat ditindaklanjuti kader.
Ketiganya menampilkan pernyataan bukan diagnosis serta versi aplikasi, model, dan stimulus.}
\end{figure}

\section*{Lampiran C. Rancangan data prospektif dan target operasional}

\textbf{Data prospektif (rancangan tahap deployment).} Rancangan kohort berpasangan di Posyandu, dengan
seluruh modalitas direkam dari anak dan sesi yang sama, acuan berupa asesmen perkembangan
terstandar dan diagnosis konsensus oleh klinisi yang dibutakan terhadap keluaran Neurogaze,
serta satu lokasi yang disisihkan sebagai validasi eksternal. Kohort ini wajib memuat anak
dengan hambatan perkembangan non-autisme sebagai kontrol negatif sulit; tanpa mereka
spesifisitas hanya terukur terhadap anak tipikal. Bagian~\ref{sec:kendala} menjelaskan
prasyarat kelembagaan yang membuka jalur ini.

\addcontentsline{toc}{section}{Lampiran C}

Seluruh nilai berikut adalah target penerimaan, bukan hasil pengukuran: alur inti berjalan
tanpa jaringan setelah aset terpasang; perangkat Android kelas menengah dengan kamera depan
minimal 5~MP pada 30~fps dan RAM minimal 4~GB; median waktu layanan 5 menit dengan P90 7
menit; kapasitas minimal 10 anak per jam per perangkat; quality gate mensyaratkan minimal
85 persen frame dengan wajah dan iris terdeteksi; presisi gaze $\leq 3^\circ$ lulus,
3-5$^\circ$ revisi, $> 5^\circ$ gagal; tidak ada unggahan media mentah pada alur MVP; setiap
hasil menyimpan versi aplikasi, model, stimulus, dan aturan.

\section*{Lampiran D. Audit dataset wajah dan keputusan karantina}
\addcontentsline{toc}{section}{Lampiran D}
\label{app:audit-wajah}

Lampiran ini melaporkan audit penuh atas dataset citra wajah statis publik yang kami tolak,
beserta angka-angka yang menjadi dasar penolakan. Tidak ada satu pun bobot atau fitur dari
eksperimen di lampiran ini yang masuk ke Neurogaze.

\textbf{Integritas teknis: lulus.} Seluruh 2.940 JPEG dapat dibaca, tidak ada berkas rusak,
tidak ada duplikat exact, dan hanya terdapat 10 pasangan near-duplicate pada pHash $\leq 4$
yang tidak satu pun melintasi label. Klasifikasi seimbang sempurna pada 1.470 citra per kelas.
Dataset ini, dengan kata lain, tidak ditolak karena kotor.

\textbf{Tata kelola: enam dari enam gagal.} Tabel~\ref{tab:governance} merangkumnya. Perlu
dibedakan lisensi platform dari persetujuan subjek: keduanya bukan hal yang sama, dan yang
kedua tidak dapat digantikan oleh yang pertama. Ketiadaan identitas partisipan berdiri
sendiri sebagai penggugur, karena tanpanya tidak ada cara membuktikan foto orang yang sama
tidak menyeberang split, sehingga estimasi performa apa pun tidak dapat ditafsirkan.

\begin{table}[h]
\centering
\caption{Metadata tata kelola yang diperiksa. Sumber: \texttt{research/hasil/audit\_wajah.json}.}
\label{tab:governance}
\small
\begin{tabular}{lll}
\hline
Metadata & Status & Akibat bila tidak ada \\
\hline
Sumber (provenance)        & tidak ada & rantai asal data tidak dapat dilacak \\
Lisensi tertulis           & tidak ada & dasar hukum pemakaian tidak jelas \\
Persetujuan subjek         & tidak ada & anak dapat diidentifikasi tanpa izin \\
Demografi                  & tidak ada & audit bias subkelompok tidak mungkin \\
Definisi label klinis      & tidak ada & arti label \emph{Autistic} tidak diketahui \\
Identitas partisipan       & tidak ada & evaluasi tidak dapat ditafsirkan \\
\hline
\end{tabular}
\end{table}

\textbf{Dua uji shortcut.} Kami menguji apakah label dapat ditebak tanpa memodelkan wajah.
Dengan lima properti berkas saja (lebar, tinggi, ukuran, rasio aspek, bytes per piksel),
regresi logistik mencapai AUC 0,4778 dan gradient boosting 0,6757. Dengan sepuluh fitur yang
menambahkan statistik piksel tingkat rendah (rerata kanal RGB, simpangan baku kecerahan,
kecerahan tepi, energi tepi), regresi logistik mencapai 0,7515 dengan permutation
$p = 0{,}005$ atas 200 permutasi. Angka kedua konsisten dengan temuan bahwa dua puluh
resolusi terbanyak seluruhnya berasal dari kelas \emph{Non\_Autistic}. Kami melaporkan
keduanya karena keduanya relevan dan arahnya berbeda: ada shortcut koleksi yang nyata,
tetapi besarnya tidak cukup menjelaskan performa CNN.

\textbf{CNN audit.} Di bawah \emph{governance gate} yang harus di-override secara sadar, CNN
EfficientNetB3 berpretrain ImageNet dilatih pada dataset ini. AUC uji mencapai 0,9324
(95\% CI 0,9070-0,9537), dengan average precision 0,9266, akurasi seimbang 0,8425, MCC
0,6861, sensitivitas 0,8143, dan spesifisitas 0,8708 pada ambang 0,5436 yang dikunci dari
validation. Split memakai grup near-duplicate dHash, yang menghilangkan kebocoran citra
identik tetapi \emph{tidak} membuat split aman terhadap identitas orang.

\textbf{Penegakan keputusan di kode.} Keputusan tidak berhenti sebagai pernyataan. Notebook
memuat enam flag tata kelola yang seluruhnya \texttt{False} dan berhenti kecuali
\texttt{OVERRIDE\_QUARANTINE} dinyatakan sadar untuk keperluan audit. Ekspor artefak dijaga
\texttt{ALLOW\_MODEL\_EXPORT = False}, sehingga sel ekspor tidak pernah berjalan dan tidak ada
bobot yang tersimpan di repositori. Dataset diberi status
\texttt{reject\_for\_asd\_modeling\_and\_quarantine}, dan \texttt{audit\_wajah.json} mencatat
syarat peninjauan ulang: provenance, lisensi, persetujuan, demografi, dan definisi label
terdokumentasi; identitas partisipan stabil yang memungkinkan evaluasi berkelompok; serta
tinjauan etik yang membenarkan pertanyaan riset yang relevan secara perilaku. Ketiganya harus
dipenuhi bersama, bukan salah satu.

\textbf{Batas penafsiran.} Angka 0,9324 tidak boleh dibaca sebagai bukti bahwa autisme dapat
dibaca dari morfologi wajah statis. Ia hanya membuktikan bahwa metrik yang terlihat meyakinkan
dapat diproduksi dari data yang tidak boleh dipakai. Artefak auditnya tersimpan pada folder
\texttt{audit\_wajah} beserta notebooknya.


\section*{Lampiran E. Tabel pendukung dan rincian hasil}

\begin{table}[h]
\caption{Hasil model pada 547 scanpath dari 54 anak, seluruhnya dengan evaluasi
dikelompokkan per anak. Baris terakhir dicantumkan sebagai pembanding representasi dan tidak
dipakai sebagai model produk.}
\label{tab:model}
\begin{center}
\small
\begin{tabular}{llrrl}
\toprule
\textbf{Model} & \textbf{Representasi} & \textbf{AUC anak} & \textbf{95\% CI} & \textbf{Status} \\
\midrule
Regresi logistik & 13 fitur geometri & 0,8228 & 0,6974-0,9290 & dibekukan sebagai MVP \\
Regresi logistik & 19 fitur penuh    & 0,8324 & 0,7069-0,9368 & butuh eye-tracker 250~Hz \\
EfficientNetB0   & citra scanpath    & 0,8819 & 0,7738-0,9681 & tidak dijalankan di produk \\
\bottomrule
\end{tabular}
\end{center}
\end{table}


\begin{table}[h]
\caption{Gate~A: hasil menurut kondisi pengujian, 100 sesi dari 25 peserta dewasa.}
\label{tab:gatea}
\begin{center}
\small
\begin{tabular}{lrrr}
\toprule
\textbf{Kondisi} & \textbf{Sesi} & \textbf{Sesi berhasil} & \textbf{Galat median} \\
\midrule
Cahaya normal   & 50 & 49 (98\%) & 1,9$^\circ$ \\
Cahaya redup    & 25 & 22 (88\%) & 2,8$^\circ$ \\
Tanpa kacamata  & 70 & 68 (97\%) & 2,0$^\circ$ \\
Dengan kacamata & 30 & 26 (87\%) & 2,9$^\circ$ \\
\midrule
Seluruh sesi    & 100 & 94 (94\%) & 2,2$^\circ$ \\
\bottomrule
\end{tabular}
\end{center}
\end{table}

\begin{table}[h]
\caption{Gate~B: distribusi AOI pada 27 pasangan siap.}
\label{tab:gateb}
\begin{center}
\small
\begin{tabular}{lrr}
\toprule
\textbf{Area perhatian} & \textbf{WebGazer} & \textbf{Neurogaze} \\
\midrule
Wajah              & 17,1628\% & 17,1374\% \\
Target kiri        & 32,7535\% & 32,8137\% \\
Target kanan       & 33,2116\% & 33,2336\% \\
Latar              & 16,8720\% & 16,8152\% \\
\bottomrule
\end{tabular}
\end{center}
\end{table}


\begin{table}[h]
\caption{Posisi Neurogaze terhadap alternatif yang tersedia. Kolom Neurogaze menyatakan
status implementasi saat makalah ini ditulis, bukan target.}
\label{tab:beda}
\begin{center}
\small
\begin{tabular}{lllll}
\toprule
\textbf{Dimensi} & \textbf{M-CHAT} & \textbf{EarliPoint} & \textbf{SenseToKnow} & \textbf{Neurogaze} \\
\midrule
Instrumen        & ceklis orang tua & eye-tracker khusus & kamera tablet & kamera tablet \\
Lokasi           & mana saja        & fasilitas klinis   & layanan primer & Posyandu \\
Operator         & orang tua        & tenaga terlatih    & tenaga klinis  & kader \\
Konektivitas     & tidak perlu      & fasilitas          & daring         & luring penuh \\
Biaya per periksa & mendekati nol   & 599 USD            & tidak dipublikasi & belum diukur; distribusi perangkat lunak statis \\
Tahap validasi   & klinis luas      & izin FDA           & prospektif 475 anak & demonstrasi rekayasa; klinis belum diuji \\
Menahan keputusan & tidak            & tidak dipublikasi  & tidak dipublikasi & ya, eksplisit \\
\bottomrule
\end{tabular}
\end{center}
\end{table}

\begin{table}[h]
\caption{Titik kerja tingkat anak pada data uji tertahan. Kolom terakhir adalah jumlah anak
yang dirujuk per 100 anak yang diskrining, yaitu besaran yang langsung menentukan beban
kapasitas rujukan.}
\label{tab:ambang}
\begin{center}
\small
\begin{tabular}{lrrrrr}
\toprule
\textbf{Ambang} & \textbf{Sensitivitas} & \textbf{Spesifisitas} & \textbf{NPP} & \textbf{NPN} & \textbf{Dirujuk/100} \\
\midrule
0,2243 (target sens.\ 0,9) & 0,885 & 0,179 & 0,500 & 0,625 & 85,2 \\
0,3500                     & 0,885 & 0,464 & 0,605 & 0,812 & 70,4 \\
0,4000                     & 0,846 & 0,607 & 0,667 & 0,810 & 61,1 \\
0,4500                     & 0,769 & 0,679 & 0,690 & 0,760 & 53,7 \\
0,4985 (Youden)            & 0,731 & 0,821 & 0,760 & 0,793 & 46,3 \\
\bottomrule
\end{tabular}
\end{center}
\end{table}

\textbf{Rincian hasil sistem dan reproduktibilitas.}

Implementasi PWA memuat Face Landmarker WASM lokal, kalibrasi, stimulus, ekstraksi fitur,
inferensi, quality gate, laporan tiga status, dan replay deterministik. Paritas numerik
antara implementasi Python (\texttt{research/features.py}) dan TypeScript
(\texttt{app/src/scanpath/features.ts}) diuji terhadap fixture bersama dan lulus pada
toleransi $10^{-12}$, baik untuk nilai fitur maupun untuk probabilitas model. Rangkaian uji
berjalan pada 16 dari 16 uji Python, 55 dari 55 uji TypeScript, dan 2 dari 2 kontrak
build/offline. Uji browser nyata memuat PWA, mengaktifkan service worker, mematikan server,
memuat ulang halaman, dan replay tetap berfungsi tanpa galat konsol.

Benchmark rantai raster $\rightarrow$ fitur $\rightarrow$ inferensi untuk 360 titik pada 250
iterasi memberi median 4,93~ms dan P90 7,67~ms pada desktop referensi (Intel Core Ultra 5
135H), yaitu margin sekitar tiga belas kali lipat terhadap ambang penerimaan 100~ms. Angka
ini kami tandai sebagai pengukuran desktop dan tidak diekstrapolasi menjadi klaim tablet;
aplikasi merekam latensi aktual setiap sesi.

\textbf{Rincian studi degradasi.} Citra sumber tidak memuat timestamp, sehingga eksperimen
30~Hz adalah sparsifikasi piksel dan tidak boleh ditafsirkan sebagai \emph{temporal
resampling}. Kondisi gabungan yang meniru perangkat sasaran adalah proxy 30~Hz ditambah
derau 1,5$^\circ$, dropout 20 persen, dan jitter pose 2$^\circ$. Jitter pose sampai
3$^\circ$ nyaris tidak berpengaruh pada kedua set; pada proxy 120~Hz set penuh masih unggul
dan pada 60~Hz keduanya bersilangan.


\begin{table}[h]
\caption{Ketahanan set fitur terhadap proxy degradasi, AUC tingkat anak. Kondisi gabungan
meniru perangkat sasaran. Seluruh angka berasal dari \texttt{research/hasil/degradasi.json}.}
\label{tab:degradasi}
\begin{center}
\small
\begin{tabular}{lrrr}
\toprule
\textbf{Kondisi} & \textbf{13 geometri} & \textbf{19 penuh} & \textbf{Selisih} \\
\midrule
Bersih (250~Hz)                       & 0,8228 & 0,8324 & $-0,0096$ \\
Derau 1,5$^\circ$                     & 0,8082 & 0,8187 & $-0,0105$ \\
Derau 3$^\circ$                       & 0,7569 & 0,7724 & $-0,0155$ \\
Dropout 30\%                          & 0,8054 & 0,8255 & $-0,0201$ \\
Jitter pose 3$^\circ$                 & 0,8214 & 0,8356 & $-0,0142$ \\
Proxy 60~Hz                           & 0,7404 & 0,7321 & $+0,0083$ \\
Proxy 30~Hz                           & 0,6882 & 0,6259 & $+0,0623$ \\
\textbf{Gabungan proxy tablet}        & \textbf{0,6827} & \textbf{0,6049} & $\mathbf{+0,0778}$ \\
\bottomrule
\end{tabular}
\end{center}
\end{table}

\begin{table}[h]
\caption{Posisi setiap tingkat validasi terhadap lingkup makalah ini. Kolom bukti hanya
memuat artefak yang ada di repositori.}
\label{tab:gate}
\begin{center}
\small
\begin{tabular}{llll}
\toprule
\textbf{Tingkat} & \textbf{Yang diuji} & \textbf{Bukti yang ada} & \textbf{Status} \\
\midrule
Algoritmik & generalisasi ke anak baru & AUC 0,8228 pada uji tertahan, 54 anak & \textbf{tuntas} \\
A & instrumen pada dewasa & 100 sesi, 25 peserta, galat median 2,2$^\circ$ & \textbf{lulus} \\
B & agreement vs WebGazer.js & 30 pasangan, 27 siap, AOI 99,7118\% & \textbf{lulus} \\
C & skrining pada balita & bukti sekunder; kohort prospektif belum ada & terbuka \\
D & operasional Posyandu & protokol tersedia; uji lapangan belum ada & terbuka \\
\bottomrule
\end{tabular}
\end{center}
\end{table}

\section*{Lampiran F. Titik kerja dan scanpath anak yang salah diklasifikasikan}

\begin{figure}[h]
\begin{center}
\includegraphics[width=0.82\linewidth]{ambang_confusion.png}
\end{center}
\caption{Titik kerja model riset Carette pada 54 anak domain sumber. (a) ambang dari target
sensitivitas 0,9; (b) titik kerja seimbang menurut indeks Youden; (c) sensitivitas,
spesifisitas, dan fraksi anak yang dirujuk sebagai fungsi ambang, dengan garis titik
menandai kedua titik kerja.}
\label{fig:ambang}
\end{figure}

\begin{figure}[h]
\begin{center}
\includegraphics[width=0.88\linewidth]{error_analysis.png}
\end{center}
\caption{Scanpath anak yang salah diklasifikasikan pada titik kerja seimbang, dengan
prediksi benar sebagai pembanding. Nilai tinta adalah fraksi piksel yang dilewati lintasan.
Kesalahan terkonsentrasi pada rekaman dengan cakupan ekstrem di kedua arah.}
\label{fig:error}
\end{figure}

\section*{Lampiran G. Rincian praproses, fitur, model, dan protokol evaluasi}
\addcontentsline{toc}{section}{Lampiran G}

Lampiran ini memuat rincian yang diringkas pada Bagian~\ref{sec:tp-ringkas}.

\begin{figure}[h]
\begin{center}
\includegraphics[width=0.82\linewidth]{pipeline_riset.png}
\end{center}
\caption{Alur penelitian Neurogaze. Pemilihan fitur, model, dan ambang seluruhnya terjadi
pada tahap 4 memakai lipatan dalam; lipatan uji tertahan hanya dibaca pada tahap 5 dan 6.
Audit dataset wajah pada Lampiran~D berjalan di luar alur ini dan tidak menyumbang bobot,
fitur, maupun keputusan apa pun ke model riset Carette.}
\label{fig:pipeline}
\end{figure}

\subsection*{Praproses sinyal pandangan}
\label{sec:tp-prepro}

Antara video mentah dan matriks fitur ada tujuh tahap, seluruhnya berjalan di perangkat
(modul \texttt{gaze/pipeline.ts}).

\begin{enumerate}
\item \textbf{Deteksi wajah dan iris.} MediaPipe Face Landmarker mengeluarkan 478 landmark
      per frame; pusat iris diproyeksikan ke kerangka koordinat lokal mata yang dibentuk
      kedua sudut kelopak, sehingga kemiringan kepala tidak terbaca sebagai perubahan arah
      pandang.
\item \textbf{Kalibrasi ke koordinat layar.} Sinyal iris $(u,v)$ dipetakan ke layar
      ternormalisasi lewat regresi affine per sumbu pada grid minimal tujuh posisi, ditutup
      koreksi drift dari target pusat. Galatnya dilaporkan dalam derajat sebagai syarat lulus
      tersendiri.
\item \textbf{Penapisan sampel.} Sampel non-finit, sampel di luar margin layar 0,12, dan
      lompatan lebih dari 0,65 satuan layar dibuang sebelum apa pun dihitung.
\item \textbf{Segmentasi deret.} Jeda di atas 180~ms, pergantian fase stimulus, atau
      pergantian epoch memutus deret. Penghalusan yang menjembatani batas ini akan
      menciptakan sakade yang tidak pernah terjadi.
\item \textbf{Penghalusan dan resampling.} Rerata bergerak eksponensial
      ($\alpha = 0{,}42$) per segmen, lalu resample 50~ms agar laju frame yang berfluktuasi
      tidak mengubah bobot tiap wilayah layar.
\item \textbf{Pemetaan area minat.} Atlas AOI membagi layar menjadi wajah, target kiri,
      target kanan, dan latar, sehingga okupansi serta latensi respons terhitung per fase.
\item \textbf{Rasterisasi.} Deret $(x,y,t)$ dirender menjadi raster $640 \times 480$
      mengikuti konvensi Carette, supaya fitur dihitung pada representasi yang sama dengan
      data pelatihan.
\end{enumerate}

\subsection*{Dua kelompok fitur scanpath}
\label{sec:tp-fitur}

Sembilan belas fitur yang dihitung dari raster terbelah menurut satu hal: apakah fitur itu
membaca warna atau tidak.

\begin{itemize}
\item \emph{Tiga belas fitur geometri} hanya membaca posisi piksel bertinta, terbagi lima
      kelompok: \emph{cakupan} (\texttt{ink\_frac}, \texttt{n\_active\_cells},
      \texttt{bbox\_fill}), \emph{dispersi} (\texttt{std\_x}, \texttt{std\_y},
      \texttt{span\_x}, \texttt{span\_y}, \texttt{aspect\_ratio}), \emph{pusat perhatian}
      (\texttt{centroid\_x}, \texttt{centroid\_y}), \emph{jarak radial}
      (\texttt{radial\_mean}, \texttt{radial\_std}), dan \emph{entropi} sebaran pada grid
      $8 \times 8$ (\texttt{grid\_entropy}).
\item \emph{Enam fitur kinematik} membaca kanal warna: rerata dan simpangan baku kecepatan
      dari kanal R, rerata akselerasi dari G, rerata jerk dari B, serta rasio fiksasi dan
      sakade dari perbandingan antar-kanal. Turunan tinggi menguatkan derau pada sampling
      rendah, jadi kelayakannya di kamera 30~fps harus diputuskan lewat pengukuran.
\end{itemize}

\subsection*{Keluarga model yang diuji}
\label{sec:tp-model}

Lima keluarga model tabular kami uji pada kedua set fitur, ditambah satu jalur citra sebagai
pembanding representasi. Semuanya memakai pemisahan berkelompok yang sama.

\begin{itemize}
\item \textbf{Regresi logistik} (ditetapkan sebagai model utama sebelum melihat hasil):
      standardisasi fitur lalu \texttt{LogisticRegression}, \texttt{max\_iter} 5000. Dipilih
      karena inferensinya cuma perkalian titik, koefisiennya dapat diekspor apa adanya ke
      \texttt{model.json}, dan kontribusi tiap fitur bisa diperiksa satu per satu.
\item \textbf{SVM kernel RBF}: batas keputusan non-linier tanpa menaikkan jumlah parameter,
      tetapi keluarannya bukan probabilitas dan bobotnya tidak terbaca manusia.
\item \textbf{Gaussian Naive Bayes}: baseline yang berguna justru karena asumsi
      independensinya jelas salah pada fitur cakupan yang saling berkorelasi.
\item \textbf{Random Forest} dan \textbf{Gradient Boosting}: ansambel pohon yang menangkap
      interaksi antar-fitur, dengan risiko menghafal pada 54 partisipan dan ukuran artefak
      jauh di atas anggaran aset luring.
\item \textbf{CNN EfficientNetB0}~\cite{tan2019} \emph{transfer learning} pada citra
      scanpath: bobot ImageNet, backbone dibekukan untuk 20 epoch pelatihan head
      (\emph{global average pooling}, dropout 0,35, satu neuron sigmoid), lalu 12 epoch
      fine-tuning pada 25 persen lapisan teratas, dengan augmentasi rotasi, translasi, zoom,
      dan kontras ringan. Bobot sampelnya \emph{equal-child}, yaitu bobot kelas dibagi jumlah
      rekaman anak, agar anak dengan 23 rekaman tidak 23 kali lebih berpengaruh.
\end{itemize}

\subsection*{Kalibrasi probabilitas dan pemilihan titik kerja}
\label{sec:tp-kalibrasi}

Skor mentah pengklasifikasi bukan probabilitas, dan alat skrining yang menampilkan angka
kepada kader wajib dikalibrasi lebih dulu.

\begin{itemize}
\item \textbf{Platt scaling} pada logit skor, dipasang hanya pada lipatan validasi, dengan
      bobot sampel invers terhadap jumlah rekaman per anak agar partisipan produktif tidak
      menyetir kalibrator.
\item \textbf{Diagnostik kalibrasi} berupa skor Brier dan kemiringan kalibrasi; kemiringan
      di atas satu menandakan probabilitas yang terlalu berhati-hati.
\item \textbf{Ambang target sensitivitas}, yaitu ambang empiris terbesar yang masih memberi
      sensitivitas minimal 0,9 pada lipatan validasi, dengan stabilitas diperiksa lewat
      bootstrap.
\item \textbf{Indeks Youden} sebagai titik kerja pembanding, dilaporkan berdampingan agar
      konsekuensi pilihan ambang terlihat.
\end{itemize}

\subsection*{Protokol evaluasi berkelompok dan uji ketahanan}
\label{sec:tp-evaluasi}

Satu anak menyumbang banyak rekaman, sehingga metode evaluasi menentukan apakah angkanya
berarti.

\begin{itemize}
\item \textbf{GroupKFold bersarang} dengan \texttt{participant\_id} sebagai grup: lima
      lipatan luar untuk uji tertahan, empat lipatan dalam untuk kalibrasi dan ambang.
\item \textbf{Agregasi ke tingkat anak} dengan merata-ratakan probabilitas rekaman, karena
      unit keputusan produk adalah anak.
\item \textbf{Bootstrap yang meresampel anak}, bukan citra, agar lebar interval kepercayaan
      mencerminkan 54 dan bukan 547.
\item \textbf{Kontrol negatif kebocoran}: pipeline yang sama dijalankan ulang dengan unit
      pengelompokan diganti, lalu penyeberangan batas split dihitung.
\item \textbf{Studi degradasi} berupa derau posisi, sparsifikasi raster sebagai proxy laju
      sampling, dropout, dan jitter pose, untuk memperkirakan performa pada perangkat sasaran
      sebelum perangkatnya tersedia.
\end{itemize}

\end{document}
